\documentclass[10pt]{article}
\usepackage[a4paper, left=2cm, right=2cm, top=2.5cm, bottom=2.5cm]{geometry}
\usepackage{fontspec}
\usepackage{xeCJK}
\usepackage{amsmath}
\usepackage{amssymb}
\usepackage{multicol}
\usepackage{xcolor}        % 须先于 minted
\usepackage[cachedir=_minted-mdtex]{minted}
\usepackage{titlesec}
\usepackage{tocloft}
\usepackage[colorlinks=true, linkcolor=black, urlcolor=blue]{hyperref}

% ===== 字体（思源宋体覆盖广，避免冷僻字方框）=====
\setCJKmainfont[BoldFont={Noto Serif CJK SC Bold}, ItalicFont={Kaiti SC}]{Noto Serif CJK SC}
\setCJKsansfont[BoldFont={Noto Sans CJK SC Bold}]{Noto Sans CJK SC}
\setCJKmonofont{Noto Sans CJK SC}

% ===== minted 全局风格 =====
\setminted{
  breaklines=true,
  breakanywhere=true,
  fontsize=\small,
  frame=single,
  bgcolor=gray!5,
  tabsize=2,
}

% ===== 常见但 LaTeX 默认未定义/易冲突的数学宏，补一下保证编译 =====
\providecommand{\abs}[1]{\left|#1\right|}
\makeatletter
\let\matrix\relax
\newcommand{\matrix}[1]{\begin{smallmatrix}#1\end{smallmatrix}}
\makeatother

% ===== 章节样式（保留 md 自带编号，故用无编号 section* + 手动入目录）=====
\setcounter{tocdepth}{2}
\titlespacing{\section}{0pt}{8pt}{4pt}
\titlespacing{\subsection}{0pt}{6pt}{3pt}
\titlespacing{\subsubsection}{0pt}{4pt}{2pt}

\setlength{\columnsep}{1cm}
\setlength{\columnseprule}{0.4pt}

\begin{document}

\begin{center}
    {\Huge\bfseries ACM 模板总结}\\[6pt]
    {\large 队友的模板（鵗曦）}
\end{center}

\tableofcontents
\vspace{1em}

\begin{multicols*}{2}

\section*{ACM模板总结}
\addcontentsline{toc}{section}{ACM模板总结}


鵗曦的板子，算法小群:1071069207 \par

注意本人的板子很多都默认包含 \par

using ll = long long; \par

const ll INF = 4e18; \par

define all(x) x.begin(), x.end() \par


\section*{杂项}
\addcontentsline{toc}{section}{杂项}


\subsection*{随机数生成器}
\addcontentsline{toc}{subsection}{随机数生成器}


\begin{minted}{cpp}
mt19937_64 rng(chrono::steady_clock::now().time_since_epoch().count());

uint64_t x = rng();                                  // 随机 64 位整数
uint64_t y = uniform_int_distribution<uint64_t>(1, (1ull<<60))(rng);  // 随机 [1, 2^60]


\end{minted}


\subsection*{重载 \_\_int128\_t 输入输出}
\addcontentsline{toc}{subsection}{重载 \_\_int128\_t 输入输出}


\begin{minted}{cpp}
ostream &operator<<(ostream &os, __int128 x) {
    if (x < 0) os << '-', x = -x;
    if (x > 9) os << x / 10;
    return os << (char) ('0' + x % 10);
}

istream &operator>>(istream &is, __int128 &x) {
    string s;
    if (!(is >> s)) return is;
    x = 0;
    size_t i = 0;
    bool neg = false;
    if (s[0] == '-') {
        neg = true;
        i = 1;
    } else if (s[0] == '+') { i = 1; }
    for (; i < s.size(); ++i) x = x * 10 + (s[i] - '0');
    if (neg) x = -x;
    return is;
}
\end{minted}




\subsection*{ACM模板}
\addcontentsline{toc}{subsection}{ACM模板}


\begin{minted}{cpp}
#include <bits/stdc++.h>
#include <ext/pb_ds/assoc_container.hpp>
#include <ext/pb_ds/tree_policy.hpp>
using namespace std;
using namespace __gnu_pbds;
using ordered_set = tree<int, null_type, less<>, rb_tree_tag, tree_order_statistics_node_update>;
using ll = long long;
using ld = long double;
using ull = unsigned long long;
constexpr const ll INF = 1e18, MOD = 1e9 + 7;
constexpr const ld EPS = 1e-7;
// #define int long long
#define all(x) x.begin(), x.end()
#ifdef LOCAL
#define  _GLIBCXX_DEBUG
#endif
#define debug1d(a) cerr << #a << " = ["; for (int i = 0; i < (int)(a).size(); i++) cerr << (i ? ", " : "") << (a)[i]; cerr << "]\n";
#define debug2d(a) cerr << #a << " = [\n"; for (int i = 0; i < (int)(a).size(); i++){ cerr << "  ["; for (int j = 0; j < (int)(a[i]).size(); j++) cerr << (j ? ", " : "") << a[i][j]; cerr << "]\n";} cerr << "]\n";
#define debug(x) cerr << "[" << #x << "] = " << (x) << "\n"
static int last_res;

int main()
{
  cin.tie(0)->sync_with_stdio(0);
  int t = 1;
  cin >> t;
  while (t--)
  {
  }
  return 0;
}
\end{minted}




\section*{数据结构}
\addcontentsline{toc}{section}{数据结构}


\subsection*{字典树}
\addcontentsline{toc}{subsection}{字典树}


\begin{minted}{cpp}
class Trie
{
    int N, idx;
    vector<vector<int>> tree;
    vector<int> cnt;
    bool sensitive;
    int alpha_size;

    int char_to_index(char ch) const
    {
        if (!sensitive) ch = tolower(ch);
        if (sensitive)
        {
            if ('a' <= ch && ch <= 'z') return ch - 'a';
            if ('A' <= ch && ch <= 'Z') return ch - 'A' + 26;
            if ('0' <= ch && ch <= '9') return ch - '0' + 52;
        }
        else
        {
            if ('a' <= ch && ch <= 'z') return ch - 'a';
            if ('0' <= ch && ch <= '9') return ch - '0' + 26;
        }
        assert(false);
        return -1;
    }

public:
    Trie(bool _sensitive = false, int _N = 3e6 + 10)
        : N(_N), idx(0), sensitive(_sensitive), cnt(N + 1, 0)
    {
        alpha_size = sensitive ? 62 : 36;
        tree.assign(N + 1, vector<int>(alpha_size, 0));
    }

    void insert(const string &s)
    {
        int point = 0;
        for (char ch : s)
        {
            int j = char_to_index(ch);
            if (!tree[point][j]) tree[point][j] = ++idx;
            point = tree[point][j];
            cnt[point]++;
        }
    }

    int query(const string &s) const
    {
        int point = 0;
        for (char ch : s)
        {
            int j = char_to_index(ch);
            if (!tree[point][j]) return 0;
            point = tree[point][j];
        }
        return cnt[point];
    }

    void clear()
    {
        for (int i = 0; i <= idx; i++)
        {
            cnt[i] = 0;
            fill(tree[i].begin(), tree[i].end(), 0);
        }
        idx = 0;
    }
};
\end{minted}




\subsection*{单调栈}
\addcontentsline{toc}{subsection}{单调栈}


\begin{minted}{cpp}
vector<int> L(n + 1), R(n + 1);
stack<int> st;

// L[i]：左边第一个严格小于 a[i] 的位置
for (int i = 1; i <= n; i++) {
  while (!st.empty() && a[st.top()] >= a[i]) st.pop();
  L[i] = st.empty() ? 0 : st.top();
  st.push(i);
}

while (!st.empty()) st.pop();

// R[i]：右边第一个严格小于 a[i] 的位置
for (int i = n; i >= 1; i--) {
  while (!st.empty() && a[st.top()] >= a[i]) st.pop();
  R[i] = st.empty() ? n + 1 : st.top();
  st.push(i);
}
\end{minted}




\subsection*{对顶堆}
\addcontentsline{toc}{subsection}{对顶堆}


\begin{minted}{cpp}

/*
  对顶堆维护中位数：
    - lo:大根堆，存较小的一半，允许比 hi 多 1 个元素
    - hi:小根堆，存较大的一半
    - 奇数个元素时，中位数 = lo.top()
*/
  priority_queue<ll> lo;                        // 大根堆（左半边）
  priority_queue<ll, vector<ll>, greater<>> hi; // 小根堆（右半边）
  for (int i = 1; i <= n; i++) {
    ll x;
    cin >> x;
    // 1. 插入到合适的堆
    if (lo.empty() || x <= lo.top())
      lo.push(x);
    else
      hi.push(x);
    // 2. 调整平衡：保证 lo.size() >= hi.size() 且差值不超过 1
    if (lo.size() > hi.size() + 1) {
      hi.push(lo.top());
      lo.pop();
    }
    if (hi.size() > lo.size()) {
      lo.push(hi.top());
      hi.pop();
    }
    // 3. 奇数个时输出中位数
    if (i & 1)
      cout << lo.top() << "\n";
  }
\end{minted}




\subsection*{并查集}
\addcontentsline{toc}{subsection}{并查集}


\begin{minted}{cpp}
//基础DUS
class DSU {
private:
  vector<int> fa, siz;
public:
  DSU(int n) {
    fa.resize(n + 1);
    siz.assign(n + 1, 1);
    iota(fa.begin(), fa.end(), 0);
  }

  int find(int x) {
    return fa[x] == x ? x : fa[x] = find(fa[x]);
  }

  bool same(int x, int y) {
    return find(x) == find(y);
  }

  bool merge(int x, int y) {
    x = find(x); y = find(y);
    if (x == y) return false;
    if (siz[x] < siz[y]) swap(x, y);
    fa[y] = x;
    siz[x] += siz[y];
    return true;
  }

  int size(int x) {
    return siz[find(x)];
  }
};

//带权DSU维护点到父节点的距离
class WeightedDSU
{
private:
  vector<int> fa, siz;
  vector<ll> dis; // dis[x] = value[x] - value[fa[x]]，find 后变成 value[x] - value[root]

public:
  WeightedDSU(int n)
  {
    fa.resize(n + 1);
    siz.assign(n + 1, 1);
    dis.assign(n + 1, 0);
    iota(fa.begin(), fa.end(), 0);
  }
  int find(int x)
  {
    if (fa[x] == x)
      return x;
    int oldFa = fa[x];
    int root = find(oldFa);
    dis[x] += dis[oldFa]; // x 到旧父亲 + 旧父亲到根 = x 到根
    fa[x] = root;
    return root;
  }
  bool same(int x, int y)
  {
    return find(x) == find(y);
  }

  // 加入约束：value[y] - value[x] = w
  // 返回 false 表示和已有关系矛盾
  bool merge(int x, int y, ll w)
  {
    int rootX = find(x);
    int rootY = find(y);

    if (rootX == rootY)
      return dis[y] - dis[x] == w;

    // 按大小合并
    if (siz[rootX] > siz[rootY])
    {
      swap(x, y);
      swap(rootX, rootY);
      w = -w;
    }
    // 让 rootX 挂到 rootY 下面
    fa[rootX] = rootY;
    siz[rootY] += siz[rootX];
    // 推导：
    // value[y] - value[x] = w
    // value[x] = dis[x] + value[rootX]
    // value[y] = dis[y] + value[rootY]
    //
    // 所以：
    // dis[y] + value[rootY] - dis[x] - value[rootX] = w
    //
    // 得：
    // value[rootX] - value[rootY] = dis[y] - dis[x] - w
    dis[rootX] = dis[y] - dis[x] - w;

    return true;
  }
  // 查询 value[x] - value[root]
  ll distToRoot(int x)
  {
    find(x);
    return dis[x];
  }
  // 查询 value[y] - value[x]
  // 调用前最好保证 same(x, y)
  ll query(int x, int y)
  {
    find(x);
    find(y);
    return dis[y] - dis[x];
  }
};
\end{minted}




\subsection*{bitset}
\addcontentsline{toc}{subsection}{bitset}


\begin{minted}{cpp}
 // 构造
  bitset<N> b(13);               // 用整数初始化，13 = 1101
  bitset<N> c(string("10101"));  // 用字符串初始化，字符串右端是第 0 位

  // 单点访问
  a[3] = 1;          // 第 3 位设为 1
  cout << a[3] << '\n';      // 访问第 3 位
  cout << a.test(3) << '\n'; // 访问第 3 位，会检查越界

  // 单点修改
  a.set(2);        // 第 2 位设为 1
  a.set(4, 1);     // 第 4 位设为 1
  a.set(4, 0);     // 第 4 位设为 0
  a.reset(2);      // 第 2 位设为 0
  a.flip(3);       // 第 3 位取反

  // 整体修改
  a.set();         // 全部设为 1
  a.reset();       // 全部设为 0
  a.flip();        // 全部取反

  // 统计和判断
  cout << a.count() << '\n'; // 1 的个数
  cout << a.size() << '\n';  // bitset 长度

  cout << a.any() << '\n';   // 是否存在至少一个 1
  cout << a.none() << '\n';  // 是否全是 0
  cout << a.all() << '\n';   // 是否全是 1

  // 位运算
  bitset<N> x(string("101010"));
  bitset<N> y(string("111000"));

  cout << (x & y) << '\n';   // 按位与
  cout << (x | y) << '\n';   // 按位或
  cout << (x ^ y) << '\n';   // 按位异或
  cout << (~x) << '\n';      // 按位取反

  // 左移和右移
  bitset<N> d(string("000101"));

  cout << (d << 2) << '\n';  // 左移 2 位
  cout << (d >> 1) << '\n';  // 右移 1 位


  // 转换
  string s = d.to_string();        // 转成 string
  unsigned long v1 = d.to_ulong(); // 转成 unsigned long
  unsigned long long v2 = d.to_ullong(); // 转成 unsigned long long

  // 输入输出
  bitset<8> inputBits;
  cin >> inputBits;          // 输入如：10110010，左边是高位
  cout << inputBits << '\n'; // 输出也是从高位到低位
\end{minted}




\subsection*{ST表}
\addcontentsline{toc}{subsection}{ST表}


\begin{minted}{cpp}
template <class T, T (*Op)(T, T)>
class ST
{
  int n, LOG;
  vector<int> lg;
  vector<vector<T>> st;
  ST() : n(0), LOG(0) {}
  void build(const vector<T> &a)
  {
    n = (int)a.size();
    lg.assign(n + 1, 0);
    for (int i = 2; i <= n; ++i)
      lg[i] = lg[i >> 1] + 1;
    LOG = lg[n] + 1;
    st.assign(LOG, vector<T>(n));
    st[0] = a;
    for (int k = 1; k < LOG; ++k)
    {
      int len = 1 << k, half = len >> 1;
      for (int i = 0; i + len <= n; ++i)
      {
        st[k][i] = Op(st[k - 1][i], st[k - 1][i + half]);
      }
    }
  }

public:
  ST(const vector<T> &a) { build(a); }
  T query(int l, int r) const
  {
    int k = lg[r - l + 1];
    return Op(st[k][l], st[k][r - (1 << k) + 1]);
  }
};
\end{minted}




\subsection*{树状数组}
\addcontentsline{toc}{subsection}{树状数组}


\begin{minted}{cpp}
// 树状数组：支持单点加 + 区间和查询
inline int lowbit(int x) { return x & -x; }
struct Fenwick {
  int n;
  vector<ll> bit;
  Fenwick(int _n) : n(_n), bit(n + 1, 0) {}


  // 单点加 delta
  void add(int idx, ll delta) {
    for (int i = idx; i <= n; i += lowbit(i))
      bit[i] += delta;
  }

  // 前缀和查询 [1..idx]
  ll sum(int idx) const {
    ll res = 0;
    for (int i = idx; i > 0; i -= lowbit(i))
      res += bit[i];
    return res;
  }

  // 区间和查询 [l..r] = sum(r) - sum(l - 1)
  ll query(int l, int r) const {
    return sum(r) - sum(l - 1);
  }
};
\end{minted}


\subsection*{Info线段树}
\addcontentsline{toc}{subsection}{Info线段树}


\begin{minted}{cpp}
struct Tag
{
  ll add = 0;

  bool empty() const
  {
    return add == 0;
  }

  void apply(const Tag &tag)
  {
    add += tag.add;
  }
};

struct Info
{
  ll sum = 0;

  void apply(const Tag &tag, int l, int r)
  {
    sum += 1LL * (r - l + 1) * tag.add;
  }
};

Info operator+(const Info &a, const Info &b)
{
  return {a.sum + b.sum};
}

template <class InfoType = Info, class TagType = Tag>
class LazySegmentTree
{
  int n;
  vector<InfoType> tree;
  vector<TagType> lazy;

  void pushUp(int id)
  {
    tree[id] = tree[id << 1] + tree[id << 1 | 1];
  }

  void apply(int id, int l, int r, const TagType &tag)
  {
    tree[id].apply(tag, l, r);
    lazy[id].apply(tag);
  }

  void pushDown(int id, int l, int r)
  {
    if (lazy[id].empty())
    {
      return;
    }

    int mid = (l + r) >> 1;
    apply(id << 1, l, mid, lazy[id]);
    apply(id << 1 | 1, mid + 1, r, lazy[id]);
    lazy[id] = TagType();
  }

  void build(int id, int l, int r, const vector<InfoType> &a)
  {
    if (l == r)
    {
      tree[id] = a[l];
      return;
    }

    int mid = (l + r) >> 1;
    build(id << 1, l, mid, a);
    build(id << 1 | 1, mid + 1, r, a);
    pushUp(id);
  }

  void modify(int id, int l, int r, int ql, int qr, const TagType &tag)
  {
    if (ql <= l && r <= qr)
    {
      apply(id, l, r, tag);
      return;
    }

    pushDown(id, l, r);

    int mid = (l + r) >> 1;

    if (ql <= mid)
    {
      modify(id << 1, l, mid, ql, qr, tag);
    }

    if (qr > mid)
    {
      modify(id << 1 | 1, mid + 1, r, ql, qr, tag);
    }

    pushUp(id);
  }

  InfoType query(int id, int l, int r, int ql, int qr)
  {
    if (ql <= l && r <= qr)
    {
      return tree[id];
    }

    pushDown(id, l, r);

    int mid = (l + r) >> 1;

    if (qr <= mid)
    {
      return query(id << 1, l, mid, ql, qr);
    }

    if (ql > mid)
    {
      return query(id << 1 | 1, mid + 1, r, ql, qr);
    }

    return query(id << 1, l, mid, ql, qr) +
           query(id << 1 | 1, mid + 1, r, ql, qr);
  }

public:
  LazySegmentTree(int n = 0)
  {
    init(n);
  }

  LazySegmentTree(const vector<InfoType> &a)
  {
    build(a);
  }

  void init(int n_)
  {
    n = n_;
    tree.assign(4 * n + 5, InfoType());
    lazy.assign(4 * n + 5, TagType());
  }

  void build(const vector<InfoType> &a)
  {
    init((int)a.size() - 1);

    if (n > 0)
    {
      build(1, 1, n, a);
    }
  }

  void modify(int l, int r, const TagType &tag)
  {
    if (l > r)
    {
      return;
    }

    modify(1, 1, n, l, r, tag);
  }

  InfoType query(int l, int r)
  {
    return query(1, 1, n, l, r);
  }
};
\end{minted}




\subsection*{主席树}
\addcontentsline{toc}{subsection}{主席树}


\begin{minted}{cpp}
// 值域可持久化线段树，空间预分配为了25*n
// 如果后续想在线插入元素，需要保证元素在建树的时候出现过,也就是提前把所有可能出现的数字都离线下来，就可以在线查询了
//空间复杂度 O(C*nlogn),单次查询、插入 O(logn)
class Segment
{
  struct node
  {
    int sum = 0;
    int leftChild = 0, rightChild = 0;
  };
  vector<node> tree;
  vector<int> root;
  vector<ll> discreted, origin;

  int idx = 0, N;

  void insert(int previousVersion, int &currentVersion, int l, int r, int position)
  {
    currentVersion = ++idx;
    tree.push_back(node());
    tree[currentVersion] = tree[previousVersion];
    tree[currentVersion].sum++;
    if (l == r)
    {
      return;
    }
    int mid = (l + r) >> 1;
    if (position <= mid)
    {
      insert(tree[previousVersion].leftChild, tree[currentVersion].leftChild, l, mid, position);
    }
    else
    {
      insert(tree[previousVersion].rightChild, tree[currentVersion].rightChild, mid + 1, r, position);
    }
  }
  ll queryKthNumber(int previousVersion, int currentVersion, int l, int r, int x)
  {
    if (l == r)
    {
      return discreted[l - 1];
    }
    int leftSum = tree[tree[currentVersion].leftChild].sum - tree[tree[previousVersion].leftChild].sum;
    int mid = (l + r) >> 1;
    if (leftSum < x)
    {
      return queryKthNumber(tree[previousVersion].rightChild, tree[currentVersion].rightChild, mid + 1, r, x - leftSum);
    }
    else
    {
      return queryKthNumber(tree[previousVersion].leftChild, tree[currentVersion].leftChild, l, mid, x);
    }
  }
  ll queryRank(int previousVersion, int currentVersion, int l, int r, int position)
  {
    if (position < l)
      return 0;

    if (r <= position)
      return tree[currentVersion].sum - tree[previousVersion].sum;

    int mid = (l + r) >> 1;
    return queryRank(tree[previousVersion].leftChild, tree[currentVersion].leftChild, l, mid, position)
     + queryRank(tree[previousVersion].rightChild, tree[currentVersion].rightChild, mid + 1, r, position);
  }
  int id(ll x) { return lower_bound(all(discreted), x) - discreted.begin() + 1; }

public:
  Segment(int n) : tree(1, node()), root(n + 1), origin(n + 1)
  {
    tree.reserve(25ll * n + 5);
    for (int i = 1; i <= n; i++)
    {
      cin >> origin[i];
      discreted.push_back(origin[i]);
    }
    sort(all(discreted));
    discreted.erase(unique(all(discreted)), discreted.end());
    N = discreted.size();
    for (int i = 1; i <= n; i++)
    {
      insert(root[i - 1], root[i], 1, N, id(origin[i]));
    }
  }

  // 查询区间第k小元素
  ll queryKthNumber(int l, int r, int k)
  {
    return queryKthNumber(root[l - 1], root[r], 1, N, k);
  }
  // 查询 X 在区间内的rank,也就是查询区间小于等于 X 的数字个数
  ll queryRank(int l, int r, ll x)
  {
    //注意x有可能未出现过，这里一定不要采用 ID
    int position = upper_bound(all(discreted),x) - discreted.begin();
    if (position == 0)
      return 0;

    return queryRank(root[l - 1], root[r], 1, N, position);
  }
};
\end{minted}


\subsection*{动态开点线段树}
\addcontentsline{toc}{subsection}{动态开点线段树}


\begin{minted}{cpp}
struct node
{
  int ch[2];
  int s;
};
// 动态开点线段树, 空间复杂度 m * 2 * log(n), m 是修改次数
class Segment {
  int cnt, n;
  // ls, rs: 左孩子,右孩子编号  sum: 区间和  lazy: 懒标记
  vector<int> ls, rs, lazy;
  vector<int> sum;

  void push_up(int u, int l, int r) {
    sum[u] = sum[ls[u]] + sum[rs[u]];
  }

  void pass_down(int u, int l, int r) {
    if (lazy[u] == -1) return;
    if (!ls[u]) ls[u] = ++cnt;
    if (!rs[u]) rs[u] = ++cnt;

    ll mid = (l + r) >> 1;
    sum[ls[u]] = lazy[u] * (mid - l + 1);
    sum[rs[u]] = lazy[u] * (r - mid);
    lazy[ls[u]] = lazy[rs[u]] = lazy[u];
    lazy[u] = -1;
  }

  // 动态开点，动态修改，注意 root 要传引用
  void modify(int &root, int l, int r, int t_l, int t_r, int k) {
    if (!root) root = ++cnt;
    if (t_l <= l && t_r >= r) {
      sum[root] = k * (r - l + 1);
      lazy[root] = k;
      return;
    }
    pass_down(root, l, r);
    ll mid = (l + r) >> 1;
    if (t_l <= mid) modify(ls[root], l, mid, t_l, t_r, k);
    if (t_r > mid) modify(rs[root], mid + 1, r, t_l, t_r, k);
    push_up(root, l, r);
  }

public:
  int root;
  explicit Segment(int n_) : root(0), cnt(0), n(n_),
    ls((ll)1e7 + 5e6, 0), rs((ll)1e7 + 5e6, 0),
    lazy((ll)1e7 + 5e6, -1), sum((ll)1e7 + 5e6, 0) {}

  void modify(int l, int r, int k) {
    modify(root, 1, n, l, r, k);
  }
};
\end{minted}


\subsection*{可持久化数组}
\addcontentsline{toc}{subsection}{可持久化数组}


\begin{minted}{cpp}
struct Node {
  int l, r;
  int val;
};
// 可持久化线段树，支持访问、修改历史版本
class Segment {
  vector<Node> tr;
  vector<int> root; // 每个版本的根节点
  int n, m, idx;

  int build(int l, int r, const vector<int> &a) {
    int x = ++idx;
    if (l == r) {
      tr[x].val = a[l];
      return x;
    }
    int mid = (l + r) >> 1;
    tr[x].l = build(l, mid, a);
    tr[x].r = build(mid + 1, r, a);
    return x;
  }

  int update(int pre, int l, int r, int pos, int val) {
    int x = ++idx;
    tr[x] = tr[pre];
    if (l == r) {
      tr[x].val = val;
      return x;
    }
    int mid = (l + r) >> 1;
    if (pos <= mid) tr[x].l = update(tr[pre].l, l, mid, pos, val);
    else tr[x].r = update(tr[pre].r, mid + 1, r, pos, val);
    return x;
  }

  int query(int x, int l, int r, int pos) {
    if (l == r) return tr[x].val;
    int mid = (l + r) >> 1;
    if (pos <= mid) return query(tr[x].l, l, mid, pos);
    else return query(tr[x].r, mid + 1, r, pos);
  }

public:
  Segment(int _n, const vector<int> &a) : n(_n), idx(0) {
    tr.resize((n + 1) * 25);  // 足够空间
    root.resize(_n + 100);    // 至少支持 n+100 操作个版本
    root[0] = build(1, n, a); // 版本 0
  }

  void modify(int i, int v, int p, int val) {
    root[i] = update(root[v], 1, n, p, val);
  }

  void clone_version(int i, int v) {
    root[i] = root[v];
  }

  int query(int v, int p) {
    return query(root[v], 1, n, p);
  }
};
\end{minted}


\subsection*{重链剖分}
\addcontentsline{toc}{subsection}{重链剖分}


\begin{minted}{cpp}
// Heavy-Light Decomposition (HLD) 重链剖分模板
class HLD {

  int n, p;
  vector<vector<int>> g;

  vector<int> siz, dep, fa;
  vector<int> son;
  vector<int> top;
  int dfs_clock;

  void dfs1(int u, int f) {
    siz[u] = 1;
    fa[u] = f;
    dep[u] = dep[f] + 1;
    int maxsz = -1;
    for (int v : g[u]) {
      if (v == f) continue;
      dfs1(v, u);
      siz[u] += siz[v];
      if (siz[v] > maxsz)
        son[u] = v, maxsz = siz[v];
    }
  }

  void dfs2(int u, int tp) {
    top[u] = tp;
    dfn[u] = ++dfs_clock;
    rnk[dfs_clock] = u;
    if (son[u]) dfs2(son[u], tp);
    for (int v : g[u]) {
      if (v == fa[u] || v == son[u]) continue;
      dfs2(v, v);
    }
    tout[u] = dfs_clock;
  }

public:
  vector<int> dfn, rnk, tout;
  void build(int root = 1) {
    dfs1(root, 0);
    dfs2(root, root);
  }
  HLD(int n, int p) : n(n), g(n + 1), p(p),
                      siz(n + 1), dep(n + 1), fa(n + 1),
                      son(n + 1), top(n + 1), dfn(n + 1), rnk(n + 1), tout(n + 1) {
    dfs_clock = 0;
  }

  void add_edge(int u, int v) {
    g[u].push_back(v);
    g[v].push_back(u);
  }

  ll query_path(int u, int v, Segment &seg) {
    ll res = 0;
    while (top[u] != top[v]) {
      if (dep[top[u]] < dep[top[v]]) swap(u, v);
      res = (res + seg.query(dfn[top[u]], dfn[u]) % p) % p;
      u = fa[top[u]];
    }
    if (dep[u] > dep[v]) swap(u, v);
    res = (res + seg.query(dfn[u], dfn[v])) % p;
    return res;
  }

  void add_path(int u, int v, int val, Segment &seg) {
    while (top[u] != top[v]) {
      if (dep[top[u]] < dep[top[v]]) swap(u, v);
      seg.modify(dfn[top[u]], dfn[u], val);
      u = fa[top[u]];
    }
    if (dep[u] > dep[v]) swap(u, v);
    seg.modify(dfn[u], dfn[v], val);
  }

  ll query_subtree(int u, Segment &seg) {
    return seg.query(dfn[u], dfn[u] + siz[u] - 1);
  }
  void modify_subtree(int u, int val, Segment &seg) {
    seg.modify(dfn[u], dfn[u] + siz[u] - 1, val);
  }

  int LCA(int u, int v) {
    while (top[u] != top[v]) {
      if (dep[top[u]] < dep[top[v]]) swap(u, v);
      u = fa[top[u]];
    }
    return dep[u] < dep[v] ? u : v;
  }
};
\end{minted}




\subsection*{线段树分治}
\addcontentsline{toc}{subsection}{线段树分治}


\begin{minted}{cpp}
auto add = [&](this auto&& self, int id, int l, int r,
               int ql, int qr, Edge e) -> void {
  if (ql <= l && r <= qr) {
    seg[id].push_back(e);
    return;
  }

  int mid = (l + r) >> 1;
  if (ql <= mid) self(id << 1, l, mid, ql, qr, e);
  if (qr > mid) self(id << 1 | 1, mid + 1, r, ql, qr, e);
};

add(1, 1, q, L, R, {u, v});
\end{minted}


\subsection*{线段树维护线性基}
\addcontentsline{toc}{subsection}{线段树维护线性基}


\begin{minted}{cpp}
// 异或线性基：维护一组数能异或出来的所有值
class LinearBasis
{
  // LOG = 60 表示最高支持 61 位整数
  static constexpr int LOG = 60;
  vector<ll> base;
  int rank = 0;
public:
  LinearBasis() : base(LOG + 1)
  {
  }
  // 清空线性基
  void clear()
  {
    fill(all(base), 0);
    rank = 0;
  }
  // 插入一个数，返回它是否增加了线性基的秩
  bool insert(ll x)
  {
    for (int bit = LOG; bit >= 0; bit--)
    {
      if ((x >> bit) & 1)
      {
        if (base[bit])
        {
          x ^= base[bit];
        }
        else
        {
          base[bit] = x;
          rank++;
          return true;
        }
      }
    }
    return false;
  }

  // 合并另一个线性基
  void merge(const LinearBasis &other)
  {
    for (int bit = LOG; bit >= 0; bit--)
    {
      if (other.base[bit])
      {
        insert(other.base[bit]);
      }
    }
  }

  // 查询能异或出的不同数的个数，包含空集
  ll querySetSize() const
  {
    return 1LL << rank;
  }
  // 查询能异或出的最大值
  ll queryMaximum() const
  {
    ll res = 0;

    for (int bit = LOG; bit >= 0; bit--)
    {
      if ((res ^ base[bit]) > res)
      {
        res ^= base[bit];
      }
    }
    return res;
  }
};

// 线段树维护区间线性基，支持单点赋值修改和区间最大异或查询
class SegmentTreeLinearBasis
{
  int n;
  vector<ll> a;
  vector<LinearBasis> tree;

  // 用左右儿子的信息更新当前节点
  void pushUp(int id)
  {
    tree[id] = tree[id << 1];
    tree[id].merge(tree[id << 1 | 1]);
  }

  // 建树，叶子节点插入单个数
  void build(int id, int l, int r)
  {
    if (l == r)
    {
      tree[id].clear();
      tree[id].insert(a[l]);
      return;
    }

    int mid = (l + r) >> 1;

    build(id << 1, l, mid);
    build(id << 1 | 1, mid + 1, r);

    pushUp(id);
  }

  // 单点赋值修改：a[pos] = val
  void modify(int id, int l, int r, int pos, ll val)
  {
    if (l == r)
    {
      a[pos] = val;
      tree[id].clear();
      tree[id].insert(val);
      return;
    }

    int mid = (l + r) >> 1;

    if (pos <= mid)
    {
      modify(id << 1, l, mid, pos, val);
    }
    else
    {
      modify(id << 1 | 1, mid + 1, r, pos, val);
    }

    pushUp(id);
  }

  // 查询区间 [ql, qr] 的线性基
  LinearBasis query(int id, int l, int r, int ql, int qr)
  {
    if (ql <= l && r <= qr)
    {
      return tree[id];
    }

    int mid = (l + r) >> 1;

    if (qr <= mid)
    {
      return query(id << 1, l, mid, ql, qr);
    }

    if (ql > mid)
    {
      return query(id << 1 | 1, mid + 1, r, ql, qr);
    }

    LinearBasis leftBasis = query(id << 1, l, mid, ql, qr);
    LinearBasis rightBasis = query(id << 1 | 1, mid + 1, r, ql, qr);

    leftBasis.merge(rightBasis);

    return leftBasis;
  }

public:
  SegmentTreeLinearBasis(int n = 0)
  {
    init(n);
  }

  SegmentTreeLinearBasis(const vector<ll> &arr)
  {
    build(arr);
  }

  // 初始化长度为 n 的线段树
  void init(int n_)
  {
    n = n_;
    a.assign(n + 1, 0);
    tree.assign(4 * n + 5, LinearBasis());
  }

  // 用 1-base 数组建树
  void build(const vector<ll> &arr)
  {
    init((int)arr.size() - 1);

    for (int i = 1; i <= n; i++)
    {
      a[i] = arr[i];
    }

    if (n > 0)
    {
      build(1, 1, n);
    }
  }

  // 单点赋值修改：a[pos] = val
  void modify(int pos, ll val)
  {
    modify(1, 1, n, pos, val);
  }

  // 查询区间 [l, r] 的线性基
  LinearBasis query(int l, int r)
  {
    return query(1, 1, n, l, r);
  }

  // 查询区间 [l, r] 能异或出的最大值
  ll queryMaximum(int l, int r)
  {
    LinearBasis res = query(l, r);
    return res.queryMaximum();
  }
};
\end{minted}




\subsection*{珂朵莉树（ODT）}
\addcontentsline{toc}{subsection}{珂朵莉树（ODT）}


\begin{minted}{cpp}
class ODT
{
    struct Node
    {
        int l, r;
        // mutable 的作用：
        // set 里的元素默认不能修改。
        // 但有些 ODT 操作会直接 it->v += x，所以 v 通常写 mutable。
        // 如果你只做区间推平，不写 mutable 也行，但建议保留。
        mutable ll v;
        bool operator<(const Node &other) const
        {
            return l < other.l;
        }
    };
    int n;
    set<Node> st;
    public:
    ODT(const vector<ll> &a)
    {
        init(a);
    }
    void init(const vector<ll> &a)
    {
        n = (int)a.size() - 1;
        st.clear();
        // 常规初始化：每个位置一个段
        for (int i = 1; i <= n; i++)
        {
            st.insert({i, i, a[i]});
        }
    }

    // split(pos) 的作用：
    // 保证存在一个左端点为 pos 的区间，并返回它的迭代器。
    // 例如当前有 [1, 10] = 5
    // split(4) 后变成：
    // [1, 3] = 5
    // [4, 10] = 5
    auto split(int pos)
    {
        // 易错点 1：
        // pos = n + 1 时，表示切在数组末尾之后。
        // 这时不需要真的插入 [n+1, ?]，直接返回 end。
        if (pos > n)
            return st.end();
        // 找到第一个左端点 > pos 的段，然后 prev 得到左端点 <= pos 的最后一个段。
        // 因为 ODT 始终覆盖 [1, n]，这个段一定包含 pos。
        auto it = prev(st.upper_bound({pos, 0, 0}));
        int L = it->l;
        int R = it->r;
        ll V = it->v;

        // 如果 pos 本来就是一个段的左端点，直接返回。
        if (L == pos)
            return it;
        // 否则把 [L, R] 切成 [L, pos - 1] 和 [pos, R]。
        st.erase(it);
        st.insert({L, pos - 1, V});

        return st.insert({pos, R, V}).first;
    }

    // 区间推平：
    // 把 [l, r] 全部赋值成 v。
    void assignRange(int l, int r, ll v)
    {
        // 易错点 2：
        // 一定要先 split(r + 1)，再 split(l)。
        // 原因：如果先 split(l)，set 结构会改变。
        auto itr = split(r + 1);
        auto itl = split(l);
        // 现在 [l, r] 内部的所有段都完整落在 [itl, itr) 里。
        st.erase(itl, itr);
        st.insert({l, r, v});
    }

    // 查询 [l, r] 的区间和。
    ll querySum(int l, int r)
    {
        auto itr = split(r + 1);
        auto itl = split(l);
        ll ans = 0;
        for (auto it = itl; it != itr; ++it)
        {
            ans += 1LL * (it->r - it->l + 1) * it->v;
        }

        return ans;
    }
};
\end{minted}


\begin{minted}{cpp}
//这里补充一个维护区间值为[D,D + 1,D + 2 ,……，D + R - L]的实际应用例子
const int V = 2 * n + 1;
    set<tuple<int, int, int>> st;
    auto split = [&](int l, int r, int d)
    {
        vector<tuple<int, int, int>> del;
        auto p = prev(st.upper_bound({l, V, V}));
        while (p != st.end())
        {
            auto [L, R, D] = *p;
            if (L > r)
                break;
            seg.add(D, D + (R - L), -1);
            del.push_back({L, R, D});
            p++;
        }
        for (auto [L, R, D] : del)
        {
            st.erase({L, R, D});
            // 左残段：[L, l - 1]
            if (L < l)
            {
                st.insert({L, l - 1, D});
                seg.add(D, D + l - 1 - L, 1);
            }

            // 右残段：[r + 1, R]
            if (R > r)
            {
                int newD = D + r + 1 - L;

                st.insert({r + 1, R, newD});
                seg.add(newD, D + R - L, 1);
            }
        }
        // 注意原段要是和旧段一样就会被erase掉，所以这里插入新段一定要在删完旧段后进行
        st.insert({l, r, d});
        seg.add(d, d + (r - l), 1);
    };
\end{minted}




\noindent\rule{\linewidth}{0.4pt}\par




\section*{数学}
\addcontentsline{toc}{section}{数学}


\subsection*{重要式子}
\addcontentsline{toc}{subsection}{重要式子}


1. $max(a,b) = \frac{a + b + \abs{a - b}}{2}$ \par

$min(a,b) = \frac{a + b - \abs{a - b}}{2}$ \par

2. $x + y  = (x\oplus y )+ 2*(x \ \&\ y)$ \par

其实就是异或负责不进位部分，按位与负责进位部分 \par

3. 变换 $ u = x + y,v = x - y$ \par

- 旧坐标系曼哈顿距离即为新坐标系的切比雪夫距离 \par

- 旧坐标系切比雪夫距离即为新坐标系的曼哈顿距离的一半 \par

- 新坐标系两点的欧几里得距离是旧坐标系的 $\sqrt2$ 倍 \par

- 切比雪夫距离：$max(|x1 - x2|, |y1 - y2|)$ \par

- 曼哈顿距离：$|u1 - u2| + |v1 - v2|$ \par

4. 斐波那契矩阵  $\matrix{10\\11}$ \par

性质：$a_{n +2}  -1= s_n$ \par

5. 整数相除精确四舍五入：$a/b = (a + b / 2)/b$ \par

6. 坐标旋转公式： \par

- 新坐标系相对原坐标系逆时针旋转 \texttt{θ}，同一个点在新坐标系下的坐标为： \par

\begin{minted}{cpp}
   X = x cosθ + y sinθ
   Y = -x sinθ + y cosθ
\end{minted}


这等价于对点做了 \texttt{-θ} 的旋转。 \par

- 点绕原点逆时针旋转 θ，点本身旋转，坐标系不动： \par

\begin{minted}{cpp}
   X = x cosθ - y sinθ
   Y = x sinθ + y cosθ
\end{minted}


7. $p$ 进制下，$\frac{a}{b}$ 能写成有限小数当且仅当化成最简分数形式之后的分母只包含p的质因子 \par

8. \par

\subsection*{公平博弈}
\addcontentsline{toc}{subsection}{公平博弈}


\subsubsection*{基本概念}
\addcontentsline{toc}{subsubsection}{基本概念}


公平组合游戏通常满足： \par

- 两人轮流操作,且\textbf{进行的操作一致}。 \par
- 游戏状态有限。 \par
- 双方\textbf{失败条件一致} \par
- 双方都采取最优策略。 \par

状态分类： \par

- 必胜态：存在一步走到必败态。 \par
- 必败态：所有后继状态都是必胜态。 \par

\noindent\rule{\linewidth}{0.4pt}\par


\subsubsection*{Bash 博弈}
\addcontentsline{toc}{subsubsection}{Bash 博弈}


模型： \par

一堆石子有 \texttt{n} 个，每次可以取 \texttt{a \textasciitilde{} b} 个，取到最后者胜。 \par

结论： \par

\begin{minted}{cpp}
n % (a + b) == 0  // 先手必败
n % (a + b) != 0  // 先手必胜
\end{minted}


\subsubsection*{Nim 博弈}
\addcontentsline{toc}{subsubsection}{Nim 博弈}


模型： \par

有多堆石子，每次选择一堆取任意正整数个，取到最后者胜。 \par

结论： \par

\begin{minted}{cpp}
xor_sum = a[1] ^ a[2] ^ ... ^ a[n];

xor_sum == 0  // 先手必败
xor_sum != 0  // 先手必胜
\end{minted}


\subsubsection*{反 Nim 博弈}
\addcontentsline{toc}{subsubsection}{反 Nim 博弈}


模型： \par

和 Nim 相同，但是取到最后一个石子的人输。 \par

结论： \par

\begin{minted}{cpp}
如果所有堆都是 1：
    堆数为奇数，先手必败
    堆数为偶数，先手必胜
否则：
    xor_sum == 0，先手必败
    xor_sum != 0，先手必胜
\end{minted}


\subsubsection*{SG 函数}
\addcontentsline{toc}{subsubsection}{SG 函数}


定义： \par

\begin{minted}{cpp}
SG(u) = mex({SG(v) | v 是 u 的后继状态})
\end{minted}


其中 \texttt{mex} 表示最小的没有出现的非负整数。 \par

结论： \par

\begin{minted}{cpp}
SG(state) == 0  // 先手必败
SG(state) != 0  // 先手必胜
\end{minted}


多个独立游戏合并： \par

\begin{minted}{cpp}
SG_total = SG1 ^ SG2 ^ ... ^ SGn;

SG_total == 0  // 先手必败
SG_total != 0  // 先手必胜
\end{minted}


\noindent\rule{\linewidth}{0.4pt}\par


\subsubsection*{SG 函数打表法}
\addcontentsline{toc}{subsubsection}{SG 函数打表法}


例如每次可以取 \texttt{moves} 中的数量： \par

\begin{minted}{cpp}
vector<int> get_sg(int N, vector<int> moves) {
    vector<int> sg(N + 1, 0);

    for (int i = 1; i <= N; i++) {
        vector<int> vis(128, 0);
        for (auto d : moves) {
            if (i >= d) {
                vis[sg[i - d]] = 1;
            }
        }

        int g = 0;
        while (vis[g]) g++;

        sg[i] = g;
    }
    return sg;
}
\end{minted}




\subsection*{快速幂}
\addcontentsline{toc}{subsection}{快速幂}


\begin{minted}{cpp}
// 快速幂：计算 a^b % mod
long long qpow(long long a, long long b, long long mod) {
  long long res = 1 % mod;
  a %= mod;
  while (b > 0) {
    if (b & 1)
      res = res * a % mod;
    a = a * a % mod;
    b >>= 1;
  }
  return res;
}

// 模意义下的乘法逆元:inv(a) = a^{-1} % mod(当 mod 为素数)
long long inv(long long a, long long mod) {
  return qpow(a, mod - 2, mod);  // 费马小定理
}
\end{minted}




\subsection*{矩阵类}
\addcontentsline{toc}{subsection}{矩阵类}


\begin{minted}{cpp}
struct Matrix {
  int n, m;
  vector<vector<ll>> a;
  ll mod;

  // 构造函数：生成 n 行 m 列、初值为 val 的矩阵，模数为 mod
  Matrix(int _n, int _m, ll val = 0, ll _mod = 1e9 + 7)
      : n(_n), m(_m), a(_n, vector<ll>(_m, val)), mod(_mod) {}

  // 构造函数：用二维 vector 初始化矩阵，自动取模
  Matrix(const vector<vector<ll>> &vec, ll _mod = 1e9 + 7) : mod(_mod) {
    n = (int)vec.size();
    m = n ? (int)vec[0].size() : 0;
    a.assign(n, vector<ll>(m));
    for (int i = 0; i < n; ++i) {
      assert((int)vec[i].size() == m);
      for (int j = 0; j < m; ++j)
        a[i][j] = (vec[i][j] % mod + mod) % mod;
    }
  }

  // 单位矩阵(继承 mod,必须是方阵)
  Matrix identity(int size) const {
    Matrix I(size, size, 0, mod);
    for (int i = 0; i < size; ++i)
      I.a[i][i] = 1;
    return I;
  }

  // 矩阵乘法（支持任意模）
  Matrix operator*(const Matrix &rhs) const {
    assert(m == rhs.n);
    Matrix res(n, rhs.m, 0, mod);
    for (int i = 0; i < n; ++i)
      for (int k = 0; k < m; ++k) {
        ll x = a[i][k] % mod;
        for (int j = 0; j < rhs.m; ++j) {
          ll y = rhs.a[k][j] % mod;
          res.a[i][j] = (res.a[i][j] + x * y % mod) % mod;
        }
      }
    return res;
  }

  // 矩阵快速幂（仅适用于方阵）
  Matrix pow(ll exp) const {
    assert(n == m);
    Matrix res = identity(n), base = *this;
    while (exp) {
      if (exp & 1) res = res * base;
      base = base * base;
      exp >>= 1;
    }
    return res;
  }

  // 打印矩阵（每行空格分隔）
  void print() const {
    for (int i = 0; i < n; ++i) {
      for (int j = 0; j < m; ++j) {
        cout << a[i][j];
        if (j + 1 < m) cout << " ";
      }
      cout << "\n";
    }
  }
};
\end{minted}


\subsection*{线性筛}
\addcontentsline{toc}{subsection}{线性筛}


\begin{minted}{cpp}
class LinearSieve
{
  int n;
  vector<bool> is_composite;
  // 可选的函数值
  bool has_phi, has_mu, has_sigma, has_d;
  // 约数和 sigma 的辅助数组：g[i] = i 最后一个素因子对应的等比和 (1 + p + ... + p^k)
  vector<ll> g;
  // 约数个数 d(n) 以及：cnt_last[i] = i 最后一个质因子 p 的指数 + 1 (即 k+1)
  vector<int> d, cnt_last;

public:
  vector<int> primes;
  vector<int> phi, mu;
  vector<ll> sigma;

  LinearSieve(int n,
              bool need_phi = false,
              bool need_mu = false,
              bool need_sigma = false,
              bool need_d = false)
      : n(n),
        is_composite(n + 1, false),
        has_phi(need_phi),
        has_mu(need_mu),
        has_sigma(need_sigma),
        has_d(need_d)
  {
    primes.reserve(n / 10);
    if (has_phi)
      phi.assign(n + 1, 0);
    if (has_mu)
      mu.assign(n + 1, 0);
    if (has_sigma)
    {
      sigma.assign(n + 1, 0);
      g.assign(n + 1, 0);
    }
    if (has_d)
    {
      d.assign(n + 1, 0);
      cnt_last.assign(n + 1, 0);
    }
    build();
  }

private:
  void build()
  {
    if (has_phi)
      phi[1] = 1;
    if (has_mu)
      mu[1] = 1;
    if (has_sigma)
    {
      sigma[1] = 1;
      g[1] = 1;
    }
    if (has_d)
    {
      d[1] = 1;
      cnt_last[1] = 1; // 随便填个 1，不会用到
    }

    for (int i = 2; i <= n; i++)
    {
      if (!is_composite[i])
      {
        primes.push_back(i);

        if (has_phi)
          phi[i] = i - 1;
        if (has_mu)
          mu[i] = -1;
        if (has_sigma)
        {
          g[i] = 1 + i;
          sigma[i] = g[i];
        }
        if (has_d)
        {
          d[i] = 2;
          cnt_last[i] = 2;
        }
      }

      for (int p : primes)
      {
        ll v = 1LL * p * i;
        if (v > n)
          break;
        int j = (int)v;
        is_composite[j] = true;

        // sigma 部分
        if (has_sigma)
        {
          if (i % p == 0)
          {
            g[j] = g[i] * p + 1;
            sigma[j] = sigma[i] / g[i] * g[j];
          }
          else
          {
            g[j] = 1 + p;
            sigma[j] = sigma[i] * g[j];
          }
        }

        // phi 部分
        if (has_phi)
        {
          if (i % p == 0)
          {
            phi[j] = phi[i] * p;
          }
          else
          {
            phi[j] = phi[i] * (p - 1);
          }
        }

        // mu 部分
        if (has_mu)
        {
          if (i % p == 0)
          {
            mu[j] = 0;
          }
          else
          {
            mu[j] = -mu[i];
          }
        }

        // d(n) 部分
        if (has_d)
        {
          if (i % p == 0)
          {
            cnt_last[j] = cnt_last[i] + 1;
            d[j] = d[i] / cnt_last[i] * cnt_last[j];
          }
          else
          {
            cnt_last[j] = 2;
            d[j] = d[i] * 2;
          }
        }
        if (i % p == 0)
          break;
      }
    }
  }
};
\end{minted}


\subsection*{质因数分解}
\addcontentsline{toc}{subsection}{质因数分解}


\begin{minted}{cpp}
const int MAXN = 1000000;

vector<int> minp(MAXN + 1), primes;

// 线性筛预处理 minp[x] = x 的最小质因子
// 分解 x 的质因数，返回 {质因子, 指数}
// 复杂度：O(质因子个数，带重数)，最坏约 O(log x)
//对1 - N 所有数字分解复杂度为O(N loglogN)
vector<pair<int, int>> factorize(int x)
{
  vector<pair<int, int>> factors;

  while (x > 1)
  {
    int p = minp[x];
    int cnt = 0;

    while (x % p == 0)
    {
      x /= p;
      cnt++;
    }

    factors.push_back({p, cnt});
  }
  return factors;
}
\end{minted}




\subsection*{扩展欧几里得}
\addcontentsline{toc}{subsection}{扩展欧几里得}


\begin{minted}{cpp}
// 裴蜀定理 Bezout Theorem
//
// 1. 对于任意整数 a, b，一定存在整数 x, y，使得：
//      a * x + b * y = gcd(a, b)
// 2. 一般地，对于多个整数 a1, a2, ..., an，
//    它们通过整数线性组合能得到的所有整数，恰好是 gcd(a1, a2, ..., an) 的倍数。
//    即方程：
//      a1 * x1 + a2 * x2 + ... + an * xn = c
//    有整数解的充要条件：
//      gcd(a1, a2, ..., an) | c


// 如何得到一组解：
//
// 1. 先用扩展欧几里得求出一组 x0, y0，使得：
//      a * x0 + b * y0 = gcd(a, b)
//
// 2. 设：
//      g = gcd(a, b)
//
//    如果要求解：
//      a * x + b * y = c
//
//    那么必须满足：
//      c % g == 0
//
// 3. 因为已经有：
//      a * x0 + b * y0 = g
//
//    两边同时乘以 c / g，可得：
//      a * (x0 * c / g) + b * (y0 * c / g) = c
//
//    所以一组特解是：
//      x = x0 * (c / g)
//      y = y0 * (c / g)
//
// 4. 通解形式：
//      x = x0 * (c / g) + (b / g) * t
//      y = y0 * (c / g) - (a / g) * t
//
//    其中 t 是任意整数。
class Exgcd {
public:
    // 扩展欧几里得：求 ax + by = gcd(a,b) 的一组解 (x,y)，返回 gcd(a,b)
    static ll extgcd(ll a, ll b, ll &x, ll &y) {
        if (b == 0) {
            x = (a >= 0 ? 1 : -1);
            y = 0;
            return llabs(a);
        }
        ll x1, y1;
        ll g = extgcd(b, a % b, x1, y1);
        // 从 b·x1 + (a%b)·y1 = g 推导
        x = y1;
        y = x1 - (a / b) * y1;
        return g;
    }

    // 模逆：求 a·inv ≡ 1 (mod m)，若 gcd(a,m)=1 返回最小非负逆元，否则返回 -1
    static ll modInverse(ll a, ll m) {
        ll x, y;
        if (extgcd(a, m, x, y) != 1) return -1;
        x %= m;
        if (x < 0) x += m;
        return x;
    }
};
\end{minted}




\subsection*{组合数(包含卡特兰数)}
\addcontentsline{toc}{subsection}{组合数(包含卡特兰数)}


\begin{minted}{cpp}
struct Combinatorics {
  int maxn;
  int mod;
  vector<long long> fac, inv, facinv;

  long long qpow(long long a, long long b) const {
    long long res = 1;
    a %= mod;
    while (b) {
      if (b & 1) res = res * a % mod;
      a = a * a % mod;
      b >>= 1;
    }
    return res;
  }

  Combinatorics(int n, int _mod) : maxn(n), mod(_mod) {
    fac.resize(n + 1);
    inv.resize(n + 1);
    facinv.resize(n + 1);
    fac[0] = facinv[0] = 1;
    for (int i = 1; i <= n; ++i)
      fac[i] = fac[i - 1] * i % mod;
    facinv[n] = qpow(fac[n], mod - 2);
    for (int i = n - 1; i >= 1; --i)
      facinv[i] = facinv[i + 1] * (i + 1) % mod;
    for (int i = 1; i <= n; ++i)
      inv[i] = fac[i] * facinv[i - 1] % mod;
  }

  long long C(int n, int k) const {
    if (k < 0 || k > n) return 0;
    return fac[n] * facinv[k] % mod * facinv[n - k] % mod;
  }

  long long A(int n, int k) const {
    if (k < 0 || k > n) return 0;
    return fac[n] * facinv[n - k] % mod;
  }
  // 卡特兰数 C_n = C(2n,n) / (n+1)
  ll catalan(int n) {
    return (C(2 * n, n) - C(2 * n, n + 1) + MOD) % MOD;
 }
  long long H(int n, int k) const {
    if (n == 0 && k == 0) return 1;
    if (n <= 0 || k < 0) return 0;
    return C(n + k - 1, k);
  }
};
\end{minted}




\subsection*{异或线性基}
\addcontentsline{toc}{subsection}{异或线性基}


\begin{minted}{cpp}
//排序采用 0-base(常识为1-base)，默认第一个数字是 0,为第 0 大，且子集是包含空集的非重集
class LinearBasis{
  //LOG = 60 表示最高支持 61 位整数
  const int LOG = 60;
  vector<ll>base,rankBase;
  int rank = 0;
  public:
  LinearBasis() : base(LOG + 1){}
  bool insert(ll x)
  {
    for(int bit = LOG ; bit >= 0 ;bit --)
    {
      if((x >> bit) & 1){
        if(base[bit]){
          x ^= base[bit];
        }else{
          base[bit] = x;
          rank ++;
          return true;
        }
      }
    }
    return false;
  }
  ll querySetSize(){
    return 1ll << rank;
  }
  ll queryMaximum(){
    ll res = 0;
    for(int bit = LOG ; bit >= 0 ;bit--){
      if((res ^ base[bit]) > res){
        res ^= base[bit];
      }
    }
    return res;
  }

  ll queryKthNumber(ll k){
      assert(k < querySetSize() && k >= 0);
    rankBase.clear();
    for(int bit = 0; bit <= LOG ; bit ++)
    {
      for(int p = bit - 1; p >=0 ; p --)
      {
        if((base[bit] >> p) & 1){
          base[bit] ^= base[p];
        }
      }
      if(base[bit]){
        rankBase.push_back(base[bit]);
      }
    }
    ll res = 0;
    for(int bit = 0 ; bit < rank ; bit ++)
    {
      if((k >> bit) & 1){
        res ^= rankBase[bit];
      }
    }
    return res;
  }
};

\end{minted}




\section*{图论}
\addcontentsline{toc}{section}{图论}


\subsection*{Dijkstra}
\addcontentsline{toc}{subsection}{Dijkstra}


\begin{minted}{cpp}
auto dijkstra = [](int s, const vector<vector<pair<ll, int>>> &adj) {
  int n = adj.size() - 1; // 1-based 节点编号,adj[0] 不使用
  ll INF = 1e18;
  vector<ll> dist(n + 1, INF);
  vector<bool> vis(n + 1, false);
  priority_queue<pair<ll, int>, vector<pair<ll, int>>, greater<>> pq;
  dist[s] = 0;
  pq.emplace(0, s);

  while (!pq.empty()) {
    auto [d, u] = pq.top();
    pq.pop();
    if (vis[u]) continue;
    vis[u] = true;
    for (auto [w, v] : adj[u]) {
      if (dist[v] > dist[u] + w) {
        dist[v] = dist[u] + w;
        pq.emplace(dist[v], v);
      }
    }
  }
  return dist; // dist[u] 表示从 s 到 u 的最短路,1-based
}
\end{minted}


\subsection*{01-BFS}
\addcontentsline{toc}{subsection}{01-BFS}


\begin{minted}{cpp}
vector<ll> zero_one_bfs(int n, int s, const vector<vector<pair<int,int>>> &g) {
    vector<ll> dist(n + 1, INF);
    deque<int> dq;

    dist[s] = 0;
    dq.emplace_back(s);

    while (!dq.empty()) {
        int u = dq.front();
        dq.pop_front();

        ll du = dist[u];
        for (auto [v, w] : g[u]) {
            if (du + w < dist[v]) {
                dist[v] = du + w;
                if (w == 0) {
                    dq.emplace_front(v); // 0 边，优先处理
                } else {
                    dq.emplace_back(v);  // 1 边，放队尾
                }
            }
        }
    }
    return dist;
}
\end{minted}




\subsection*{拓扑排序}
\addcontentsline{toc}{subsection}{拓扑排序}


\begin{minted}{cpp}
  while (!q.empty())
  {
    int sz = (int)q.size();
    vector<int> layer;
    layer.reserve(sz);
    // 这一轮队列里的点就是当前这一层
    while (sz--)
    {
      int u = q.front();
      q.pop();
      layer.push_back(u);
      for (int v : g[u])
      {
        if (--indeg[v] == 0)
        {
          q.push(v);
          // level[v] = (level[u] == -1 ? 0 : level[u] + 1); // 如果想标层号，可以改
        }
      }
    }
    layers.push_back(move(layer));
  }
\end{minted}




\subsection*{匈牙利二分图最大匹配}
\addcontentsline{toc}{subsection}{匈牙利二分图最大匹配}


\begin{minted}{cpp}
/*
  Hungarian (DFS 版) — 二分图最大匹配
  顶点集:L = {1..nL}, R = {1..nR}   —— 1-based
  用法：
    Hungarian g(nL, nR);
    g.add_edge(u, v);                // u∈L, v∈R
    int ans = g.max_matching();      // 返回最大匹配数
    auto M = g.get_matching();       // 返回所有匹配对 (u,v)
  说明：
    - 多重边自动等价（不会影响结果）。
    - 每次 max_matching() 会在内部从空匹配开始跑一遍；若要增量维护需自己管理。
*/
struct Hungarian {
  int nL, nR;
  vector<vector<int>> adj;     // L侧邻接:adj[u] = {v1, v2, ...} , v∈R
  vector<int> matchL, matchR;  // matchL[u] = u 匹配到的 R点:matchR[v] = v 匹配到的 L点
  vector<int> vis;             // DFS 访问标记（在 R 侧）

  Hungarian(int nL = 0, int nR = 0) { init(nL, nR); }

  void init(int _nL, int _nR) {
    nL = _nL; nR = _nR;
    adj.assign(nL + 1, {});
    matchL.assign(nL + 1, 0);
    matchR.assign(nR + 1, 0);
    vis.assign(nR + 1, 0);
  }
  inline void add_edge(int u, int v) {
    // 要求:1 <= u <= nL, 1 <= v <= nR
    adj[u].push_back(v);
  }
  // 尝试从 L侧点 u 增广
  bool dfs(int u) {
    for (int v : adj[u]) {
      if (vis[v]) continue;
      vis[v] = 1;
      // 若 v 未匹配，或能将 v 当前匹配的左点“让位”出去
      if (matchR[v] == 0 || dfs(matchR[v])) {
        matchL[u] = v;
        matchR[v] = u;
        return true;
      }
    }
    return false;
  }
  // 主过程：返回最大匹配数
  int max_matching() {
    // 清空旧匹配（确保可重复调用）
    fill(matchL.begin(), matchL.end(), 0);
    fill(matchR.begin(), matchR.end(), 0);
    int mcbm = 0;
    // 朴素匈牙利：对每个未匹配的左点，清空 vis，跑一遍 DFS 增广
    for (int u = 1; u <= nL; ++u) {
      if (matchL[u]) continue;
      fill(vis.begin(), vis.end(), 0);
      if (dfs(u)) ++mcbm;
    }
    return mcbm;
  }
  // 返回所有匹配对 (u,v)
  vector<pair<int,int>> get_matching() const {
    vector<pair<int,int>> res;
    res.reserve(nL);
    for (int u = 1; u <= nL; ++u) {
      if (matchL[u]) res.emplace_back(u, matchL[u]);
    }
    return res;
  }
};
\end{minted}


\subsection*{欧拉回路}
\addcontentsline{toc}{subsection}{欧拉回路}


\begin{minted}{cpp}
// 无向图欧拉回路 demo：Hierholzer 算法
//
// 结论：
// 1. 无向图存在欧拉回路的必要条件：所有点度数都是偶数。
// 2. 如果题目要求“整张有边图形成一条欧拉回路”，还需要所有有边的点连通。
// 3. 如果题目允许每个连通块分别处理，则不需要整图连通；只要每个点度数为偶数，每个有边连通块都有欧拉回路。
//
// 实现易错点：
// 1. 无向边在邻接表中存两次，但两条反向边必须共用同一个 id。
// 2. used[id] 表示原始边是否已经走过，大小应该按边数 m 开。
// 3. ptr[u] 记录 u 的邻接表扫描到哪里，避免每次从头扫导致复杂度退化。
// 4. ptr[u]++ 必须写在判断 used[id] 前面，否则遇到已用边可能死循环。
// 5. 如果要输出欧拉回路顺序，通常在递归返回后 push 边，最后 reverse。
// 6. 如果只是在欧拉遍历过程中给边定向 / 统计，不一定需要保存完整路径。
struct Edge
{
    int v, id;
};
void solve()
{
    int n, m;
    cin >> n >> m;
    vector<vector<Edge>> graph(n + 1);
    vector<int> degree(n + 1, 0);
    for (int id = 1; id <= m; id++)
    {
        int u, v;
        cin >> u >> v;
        // 无向边存两次，但共用同一个 id。
        graph[u].push_back({v, id});
        graph[v].push_back({u, id});
        degree[u]++;
        degree[v]++;
    }
    // 欧拉回路要求每个点度数为偶数。
    for (int i = 1; i <= n; i++)
    {
        if (degree[i] & 1)
        {
            cout << "-1\n";
            return;
        }
    }
    vector<int> ptr(n + 1, 0);
    vector<char> used(m + 1, false);
    vector<int> eulerEdges; // 欧拉回路的边序
    vector<int> eulerNodes; // 欧拉回路的点序
    auto dfs = [&](this auto &&self, int u) -> void
    {
        while (ptr[u] < (int)graph[u].size())
        {
            auto [v, id] = graph[u][ptr[u]];
            // 必须先移动 ptr[u]，否则遇到 used[id] 时会卡在原地。
            ptr[u]++;
            // 无向边在邻接表里出现两次，用 used[id] 防止重复走。
            if (used[id])
                continue;
            used[id] = true;
            self(v);
            // Hierholzer 的路径拼接在回溯阶段完成。
            // 所以这里回溯后记录边，最后 reverse。
            eulerEdges.push_back(id);
        }
        // 记录点序也是回溯时加入，最后 reverse。
        eulerNodes.push_back(u);
    };
    int start = 1;
    for (int i = 1; i <= n; i++)
    {
        if (degree[i])
        {
            start = i;
            break;
        }
    }
    dfs(start);
    reverse(eulerEdges.begin(), eulerEdges.end());
    reverse(eulerNodes.begin(), eulerNodes.end());
    // 如果是整张有边图一条欧拉回路，应当刚好走完 m 条边。
    if ((int)eulerEdges.size() != m)
    {
        cout << "-1\n";
        return;
    }
}

\end{minted}




\subsection*{树上差分}
\addcontentsline{toc}{subsection}{树上差分}


\begin{minted}{cpp}
// 树上差分模板整理
// 1. 路径边差分：统计每条边被多少条路径经过
// 2. 路径点差分：统计每个点被多少条路径经过
// 3. DFS 序子树差分：子树加、补子树加，Meet 这题用的是这一类
//
// 默认已有：
// g[u]           : 树的邻接表
// up[0][u]       : u 的父亲，根的父亲为自身
// tin[u], tout[u]: DFS 序
// lca(u, v)      : 求 LCA
// n              : 点数

// ============================================================
// 1. 路径边差分
// 给路径 u-v 上所有边加 1
// 回收后 diff[x] 表示边 x - up[0][x] 被经过多少次，根节点无意义
// ============================================================

vector<int> edgeDiff(n + 1, 0);
vector<int> edgeCnt(n + 1, 0);

auto addPathEdge = [&](int u, int v)
{
    int w = lca(u, v);

    edgeDiff[u]++;
    edgeDiff[v]++;
    edgeDiff[w] -= 2;
};

auto collectEdge = [&](this auto &&self, int u, int fa) -> void
{
    for (int v : g[u])
    {
        if (v == fa)
            continue;

        self(v, u);

        edgeDiff[u] += edgeDiff[v];
    }

    edgeCnt[u] = edgeDiff[u];
};

// 使用：
// for (auto [u, v] : paths)
//     addPathEdge(u, v);
// collectEdge(1, 1);


// ============================================================
// 2. 路径点差分
// 给路径 u-v 上所有点加 1
// 回收后 pointDiff[x] 表示点 x 被多少条路径经过
// ============================================================

vector<int> pointDiff(n + 1, 0);

auto addPathPoint = [&](int u, int v)
{
    int w = lca(u, v);

    pointDiff[u]++;
    pointDiff[v]++;
    pointDiff[w]--;

    if (up[0][w] != w)
        pointDiff[up[0][w]]--;
};

auto collectPoint = [&](this auto &&self, int u, int fa) -> void
{
    for (int v : g[u])
    {
        if (v == fa)
            continue;

        self(v, u);

        pointDiff[u] += pointDiff[v];
    }
};

// 使用：
// for (auto [u, v] : paths)
//     addPathPoint(u, v);
// collectPoint(1, 1);


// ============================================================
// 3. DFS 序子树差分 / 补子树差分
// 适合：子树整体加、整棵树减去某子树整体加
// Meet 这题的禁止连通块标记用的是这一类
// ============================================================

vector<int> diff(n + 3, 0);
int globalAdd = 0;

auto addSubtree = [&](int u, int val)
{
    diff[tin[u]] += val;
    diff[tout[u] + 1] -= val;
};

auto addAllExceptSubtree = [&](int u, int val)
{
    globalAdd += val;
    addSubtree(u, -val);
};

// 切断边 a-b，给包含 a 的那一侧加 1
auto addSide = [&](int a, int b)
{
    if (up[0][a] == b)
    {
        // a 是 b 的儿子，包含 a 的一侧就是 subtree(a)
        addSubtree(a, 1);
    }
    else
    {
        // b 是 a 的儿子，包含 a 的一侧就是整棵树 - subtree(b)
        addAllExceptSubtree(b, 1);
    }
};

// 按 DFS 序回收贡献
// cur 表示当前 DFS 序位置对应点的贡献值
auto recoverDfnDiff = [&]()
{
    vector<int> value(n + 1, 0);

    int cur = globalAdd;
    for (int i = 1; i <= n; i++)
    {
        cur += diff[i];
        value[i] = cur;
    }

    return value;
};

// 使用：
// addSubtree(u, 1);              // 子树 u 加 1
// addAllExceptSubtree(u, 1);     // 整棵树 - subtree(u) 加 1
// addSide(a, b);                 // 切断 a-b 后，包含 a 的一侧加 1
// auto value = recoverDfnDiff(); // value[i] 是 DFS 序位置 i 对应点的贡献
\end{minted}




\subsection*{树上启发式合并}
\addcontentsline{toc}{subsection}{树上启发式合并}


\begin{minted}{cpp}
vector<int> sz(n + 1), son(n + 1);
function<void(int, int)> dfs0 = [&](int u, int parent) {
  sz[u] = 1;
  ll maxi = 0;
  for (int v : adj[u]) {
    if (v == parent) continue;
    dfs0(v, u);
    if (sz[v] > maxi) {
      maxi = sz[v];
      son[u] = v;
    }
    sz[u] += sz[v];
  }
};
dfs0(1, 0);

function<void(int, int, int)> add = [&](int u, int parent, int child) {
  for (int v : adj[u]) {
    if (v == parent || v == child) continue;
    // TODO: implement add logic
  }
};

function<void(int, int)> clean = [&](int u, int parent) {
  for (int v : adj[u]) {
    if (v == parent) continue;
    // TODO: implement clean logic
  }
};

function<void(int, int, bool)> dfs = [&](int u, int parent, bool keep) {
  for (int v : adj[u]) {
    if (v == parent || v == son[u]) continue;
    dfs(v, u, false);
  }
  if (son[u]) dfs(son[u], u, true);
  add(u, parent, son[u]);

  if (!keep) {
    clean(u, parent);
  }
};
dfs(1, 0, true);
\end{minted}


\subsection*{倍增求LCA}
\addcontentsline{toc}{subsection}{倍增求LCA}


\begin{minted}{cpp}
// 倍增求 LCA + 距离查询
struct LCA {
  int n, LOG;
  vector<vector<int>> up, adj;
  vector<int> depth;

  LCA(int _n) : n(_n) {
    LOG = 32 - __builtin_clz(n);
    up.assign(n + 1, vector<int>(LOG));
    adj.resize(n + 1);
    depth.resize(n + 1);
  }

  void add_edge(int u, int v) {
    adj[u].push_back(v);
    adj[v].push_back(u);
  }

  void dfs(int u, int p) {
    up[u][0] = p;
    for (int i = 1; i < LOG; ++i)
      up[u][i] = up[up[u][i - 1]][i - 1];
    for (int v : adj[u]) {
      if (v != p) {
        depth[v] = depth[u] + 1;
        dfs(v, u);
      }
    }
  }

  void init(int root = 1) {
    depth[root] = 0;
    dfs(root, 0);
  }

  int lca(int u, int v) const {
    if (depth[u] < depth[v]) swap(u, v);
    for (int i = LOG - 1; i >= 0; --i) {
      if (depth[u] - (1 << i) >= depth[v])
        u = up[u][i];
    }
    if (u == v) return u;
    for (int i = LOG - 1; i >= 0; --i) {
      if (up[u][i] != up[v][i]) {
        u = up[u][i];
        v = up[v][i];
      }
    }
    return up[u][0];
  }

  int get_distance(int u, int v) const {
    int _lca = lca(u, v);
    return depth[u] + depth[v] - 2 * depth[_lca];
  }
  int get_kth_ancestor(int u, int k) {
    for (int j = LOG; j >= 0; j--) {
        if (k & (1 << j)) {
            u = up[u][j];
            if (u == 0) break;
        }
    }
    return u;
  }
};
\end{minted}


\subsection*{Tarjan\_SCC}
\addcontentsline{toc}{subsection}{Tarjan\_SCC}


\begin{minted}{cpp}
class TarjanSCC
{
public:
  int n, scc_cnt = 0;
  vector<vector<int>> adj;            // 原图邻接表
  vector<int> dfn, low, stk, scc_id;  // Tarjan 必需数组
  vector<bool> in_stk;
  vector<vector<int>> sccs;           // 强连通分量集合
  int timestamp = 0, top = 0;

  TarjanSCC(int _n) : n(_n),
    adj(_n + 1), dfn(_n + 1), low(_n + 1),
    stk(_n + 1), scc_id(_n + 1), in_stk(_n + 1, false) {}

  /// 添加有向边
  void add_edge(int u, int v) {
    adj[u].emplace_back(v);
  }

  /// 运行 Tarjan 算法，提取所有 SCC
  void run() {
    for (int u = 1; u <= n; ++u) {
      if (!dfn[u]) dfs(u);
    }
  }

  /// 构建缩点图(有向无环图 DAG),去掉重边
  vector<vector<int>> build_dag() {
    vector<vector<int>> dag(scc_cnt + 1);
    vector<unordered_set<int>> dag_set(scc_cnt + 1); // 去重用
    for (int u = 1; u <= n; ++u) {
      for (int v : adj[u]) {
        int su = scc_id[u], sv = scc_id[v];
        if (su != sv && !dag_set[su].count(sv)) {
          dag[su].push_back(sv);
          dag_set[su].insert(sv);
        }
      }
    }
    return dag;
  }

private:
  void dfs(int u) {
    dfn[u] = low[u] = ++timestamp;
    stk[++top] = u; in_stk[u] = true;
    for (int v : adj[u]) {
      if (!dfn[v]) {
        dfs(v);
        low[u] = min(low[u], low[v]);
      } else if (in_stk[v]) {
        low[u] = min(low[u], dfn[v]);
      }
    }
    if (dfn[u] == low[u]) {
      ++scc_cnt;
      sccs.emplace_back();
      int x;
      do {
        x = stk[top--];
        in_stk[x] = false;
        scc_id[x] = scc_cnt;
        sccs.back().push_back(x);
      } while (x != u);
    }
  }
};

\end{minted}


\subsection*{Tarjan\_EBCC}
\addcontentsline{toc}{subsection}{Tarjan\_EBCC}


\begin{minted}{cpp}
// Tarjan_EBCC: 在原图上(带边编号)求割边(桥)并缩点为边双连通分量(edge-biconnected components)
class Tarjan_EBCC {
    int n;
    vector<int> dfn, low;
    int timer;
    vector<int> bridges;
    vector<int> comp;
    int comp_cnt;

public:
    vector<vector<pair<int,int>>> g;
    Tarjan_EBCC(int n = 0) : n(n), g(n+1), dfn(n+1,0), low(n+1,0), timer(0), comp(n+1,0), comp_cnt(0) {}

    void add_edge(int u, int v, int id) {
        g[u].push_back({v,id});
        g[v].push_back({u,id});
    }

    void run() {
        for(int i=1;i<=n;i++){
            if(!dfn[i]) dfs_bridge(i,-1);
        }
        unordered_set<int> bridgeSet;
        bridgeSet.reserve(bridges.size()*2+10);
        for(int eid: bridges) bridgeSet.insert(eid);
        for(int i=1;i<=n;i++){
            if(!comp[i]){
                ++comp_cnt;
                dfs_comp(i, comp_cnt, bridgeSet);
            }
        }
    }

    vector<int> get_bridges() const { return bridges; }
    vector<int> get_components() const { return comp; }
    int get_comp_cnt() const { return comp_cnt; }

private:
    void dfs_bridge(int u, int peid) {
        dfn[u] = low[u] = ++timer;
        for(auto [v,eid] : g[u]){
            if(!dfn[v]){
                dfs_bridge(v,eid);
                low[u] = min(low[u], low[v]);
                if(low[v] > dfn[u]) bridges.push_back(eid);
            } else if(eid != peid){
                low[u] = min(low[u], dfn[v]);
            }
        }
    }

    void dfs_comp(int u, int cid, const unordered_set<int> &bridgeSet) {
        comp[u] = cid;
        for(auto [v,eid] : g[u]){
            if(bridgeSet.count(eid)) continue;
            if(!comp[v]) dfs_comp(v, cid, bridgeSet);
        }
    }
};
\end{minted}


\subsection*{Tarjan\_VBCC}
\addcontentsline{toc}{subsection}{Tarjan\_VBCC}


\begin{minted}{cpp}
class Tarjan_VBCC
{
  struct Edge { int to, id; };

  int n, eid;
  vector<vector<Edge>> g;
  vector<int> dfn, low, is_cut;
  int timer;
  vector<pair<int,int>> estk;
  vector<vector<int>> bcc;
  unordered_set<long long> edge_set;

  static long long encode(int u, int v) {
    if (u > v) swap(u, v);
    return ((long long)u << 32) | (unsigned)v;
  }

public:
  explicit Tarjan_VBCC(int n = 0) { init(n); }

  /// 初始化
  void init(int N) {
    n = N; eid = 0;
    g.assign(n + 1, {});
    dfn.assign(n + 1, 0);
    low.assign(n + 1, 0);
    is_cut.assign(n + 1, 0);
    timer = 0;
    estk.clear(); bcc.clear(); edge_set.clear();
  }

  /// 添加无向边，自动去掉自环和重边
  void add_edge(int u, int v) {
    if (u == v) return;
    long long code = encode(u, v);
    if (edge_set.count(code)) return;
    edge_set.insert(code);
    ++eid;
    g[u].push_back({v, eid});
    g[v].push_back({u, eid});
  }

private:
  /// Tarjan 深搜，维护 dfn/low，识别点双
  void dfs(int u, int parent_eid) {
    dfn[u] = low[u] = ++timer;
    int child = 0;
    for (auto [v, id] : g[u]) {
      if (!dfn[v]) {
        estk.emplace_back(u, v);
        ++child;
        dfs(v, id);
        low[u] = min(low[u], low[v]);
        if (low[v] >= dfn[u]) {
          if (parent_eid != -1 || child > 1) is_cut[u] = 1;
          unordered_set<int> nodes;
          while (!estk.empty()) {
            auto e = estk.back(); estk.pop_back();
            nodes.insert(e.first); nodes.insert(e.second);
            if (e.first == u && e.second == v) break;
          }
          vector<int> comp(nodes.begin(), nodes.end());
          bcc.push_back(move(comp));
        }
      } else if (id != parent_eid && dfn[v] < dfn[u]) {
        estk.emplace_back(u, v);
        low[u] = min(low[u], dfn[v]);
      }
    }
  }

public:
  /// 运行 Tarjan 算法，提取所有点双分量
  void run() {
    for (int i = 1; i <= n; ++i) {
      if (!dfn[i]) {
        dfs(i, -1);
        if (!estk.empty()) {
          unordered_set<int> nodes;
          while (!estk.empty()) {
            auto e = estk.back(); estk.pop_back();
            nodes.insert(e.first); nodes.insert(e.second);
          }
          vector<int> comp(nodes.begin(), nodes.end());
          bcc.push_back(move(comp));
        }
      }
    }
  }

  /// 返回所有点双分量
  const vector<vector<int>>& get_bcc() const { return bcc; }

  /// 返回所有割点
  vector<int> get_cutpoints() const {
    vector<int> res;
    for (int i = 1; i <= n; ++i)
      if (is_cut[i]) res.push_back(i);
    return res;
  }
};

\end{minted}


\subsection*{最小割最大流}
\addcontentsline{toc}{subsection}{最小割最大流}


\begin{minted}{cpp}
一般图最大流总复杂度：
O(n^2 m)
二分图匹配、单位容量网络等特殊图上通常更快，常见复杂度：
O(m sqrt(n))
  
 class Dinic
{
private:
  struct Edge
  {
    int from, to;
    ll cap;       // 当前剩余容量
    ll originCap; // 原始容量，用于求割边代价、查询流量
    int id;       // 输入边编号，不需要输出割边时可以不管
  };

  int n;
  vector<Edge> edges;
  vector<vector<int>> g;
  vector<int> level, cur;
  vector<int> edgeIdToIndex;

  bool bfs(int s, int t)
  {
    fill(all(level), -1);

    queue<int> q;
    level[s] = 0;
    q.push(s);

    while (!q.empty())
    {
      int u = q.front();
      q.pop();

      for (int edgeIndex : g[u])
      {
        Edge &e = edges[edgeIndex];

        if (e.cap > 0 && level[e.to] == -1)
        {
          level[e.to] = level[u] + 1;
          q.push(e.to);
        }
      }
    }

    return level[t] != -1;
  }

  ll dfs(int u, int t, ll pushed)
  {
    if (u == t)
      return pushed;
    if (pushed == 0)
      return 0;

    for (int &i = cur[u]; i < (int)g[u].size(); i++)
    {
      int edgeIndex = g[u][i];
      Edge &e = edges[edgeIndex];

      if (e.cap == 0)
        continue;
      if (level[e.to] != level[u] + 1)
        continue;

      ll tr = dfs(e.to, t, min(pushed, e.cap));
      if (tr == 0)
        continue;

      e.cap -= tr;
      edges[edgeIndex ^ 1].cap += tr;

      return tr;
    }

    return 0;
  }

public:
  // 建立一个 1-base 流图。
  // 可用点编号：1, 2, ..., n。
  // 0 号点默认空着不用。
  Dinic(int n = 0)
  {
    init(n);
  }

  // 重新初始化图。
  // 多测时可以用 dinic.init(n) 清空旧图。
  void init(int n_)
  {
    n = n_;
    edges.clear();
    g.assign(n + 1, {});
    level.assign(n + 1, -1);
    cur.assign(n + 1, 0);
    edgeIdToIndex.clear();
  }

  // 加一条有向边 from -> to，容量为 cap。
  //
  // id 是这条边的外部编号。
  // 如果你只要求最大流，不需要传 id。
  // 如果你要输出最小割边，建议传入输入边编号 i。
  //
  // 例：
  // dinic.addEdge(u, v, c);
  // dinic.addEdge(u, v, c, i);
  void addEdge(int from, int to, ll cap, int id = -1)
  {
    if (id >= 0)
    {
      if ((int)edgeIdToIndex.size() <= id)
      {
        edgeIdToIndex.resize(id + 1, -1);
      }
      edgeIdToIndex[id] = (int)edges.size();
    }

    edges.push_back({from, to, cap, cap, id});
    g[from].push_back((int)edges.size() - 1);

    edges.push_back({to, from, 0, 0, -1});
    g[to].push_back((int)edges.size() - 1);
  }

  // 加一条无向边 u - v，两个方向容量都是 cap。
  //
  // 注意：
  // 这是“无向容量边”，等价于：
  // u -> v 容量 cap
  // v -> u 容量 cap
  //
  // 如果题目是普通无向图最小割，通常可以这样建。
  void addUndirectedEdge(int u, int v, ll cap, int id = -1)
  {
    addEdge(u, v, cap, id);
    addEdge(v, u, cap, id);
  }

  // 求 s 到 t 的最大流。
  //
  // 返回值：
  // 从 s 最多能流到 t 的流量。
  //
  // 根据最大流最小割定理，这个值也等于 s-t 最小割容量。
  ll maxFlow(int s, int t)
  {
    ll flow = 0;

    while (bfs(s, t))
    {
      fill(all(cur), 0);

      while (true)
      {
        ll pushed = dfs(s, t, INF);
        if (pushed == 0)
          break;
        flow += pushed;
      }
    }

    return flow;
  }

  // 跑完 maxFlow(s, t) 之后调用。
  //
  // 返回一个 vis 数组：
  // vis[u] = 1 表示 u 在最小割的 S 侧。
  // vis[u] = 0 表示 u 在最小割的 T 侧。
  //
  // 原理：
  // 在残量网络里，从 s 出发，只走剩余容量 > 0 的边。
  vector<int> minCutSide(int s)
  {
    vector<int> vis(n + 1, 0);
    queue<int> q;

    vis[s] = 1;
    q.push(s);

    while (!q.empty())
    {
      int u = q.front();
      q.pop();

      for (int edgeIndex : g[u])
      {
        Edge &e = edges[edgeIndex];

        if (e.cap > 0 && !vis[e.to])
        {
          vis[e.to] = 1;
          q.push(e.to);
        }
      }
    }

    return vis;
  }

  // 跑完 maxFlow(s, t) 之后调用。
  //
  // 返回一组最小割边的输入编号。
  //
  // 一条原图边 u -> v 属于这组最小割，当且仅当：
  // u 在 S 侧，v 在 T 侧。
  //
  // 使用这个函数时，addEdge 最好传入 id：
  // dinic.addEdge(u, v, c, i);
  vector<int> minCutEdges(int s)
  {
    vector<int> vis = minCutSide(s);
    vector<int> cutEdges;

    for (int edgeIndex = 0; edgeIndex < (int)edges.size(); edgeIndex += 2)
    {
      Edge &e = edges[edgeIndex];

      if (e.id != -1 && vis[e.from] && !vis[e.to])
      {
        cutEdges.push_back(e.id);
      }
    }

    return cutEdges;
  }

  // 跑完 maxFlow(s, t) 之后调用。
  //
  // 返回当前这组最小割边的原始容量之和。
  //
  // 正常情况下：
  // dinic.minCutCost(s) == dinic.maxFlow(s, t)
  ll minCutCost(int s)
  {
    vector<int> vis = minCutSide(s);
    ll cost = 0;

    for (int edgeIndex = 0; edgeIndex < (int)edges.size(); edgeIndex += 2)
    {
      Edge &e = edges[edgeIndex];

      if (vis[e.from] && !vis[e.to])
      {
        cost += e.originCap;
      }
    }

    return cost;
  }

  // 查询某条输入边实际流了多少。
  //
  // 前提：
  // addEdge 时传入过这个输入边编号 id。
  //
  // 例：
  // dinic.addEdge(u, v, c, i);
  // ...
  // ll used = dinic.getFlow(i);
  ll getFlow(int id)
  {
    int edgeIndex = edgeIdToIndex[id];
    return edges[edgeIndex].originCap - edges[edgeIndex].cap;
  }
};

\end{minted}


\subsection*{最小费用最大流}
\addcontentsline{toc}{subsection}{最小费用最大流}


\begin{minted}{cpp}
// 最小费用最大流
// useSPFA = false：适用于原始边费用非负
// useSPFA = true ：适用于原始边费用可能为负
// 复杂度：O(SPFA + A * m log n)，A 为增广次数
class MinCostMaxFlow {
private:
    struct Edge {
        int from, to; // 边的起点和终点
        ll cap;       // 当前残余容量
        ll cost;      // 单位流量费用
    };
    int n;
    vector<Edge> edges;
    vector<vector<int>> g;
    // SPFA 初始化势能 h
    // 作用：当原始费用可能为负时，先算出从 s 到每个点的最短路势能
    void initPotentialBySPFA(int s, vector<ll> &h) {
        vector<int> inq(n + 1, 0);
        queue<int> q;
        fill(all(h), INF);
        h[s] = 0;
        q.push(s);
        inq[s] = 1;
        while (!q.empty()) {
            int u = q.front();
            q.pop();
            inq[u] = 0;
            for (int id : g[u]) {
                Edge &e = edges[id];
                if (e.cap <= 0) continue;
                int v = e.to;
                if (h[v] > h[u] + e.cost) {
                    h[v] = h[u] + e.cost;

                    if (!inq[v]) {
                        inq[v] = 1;
                        q.push(v);
                    }
                }
            }
        }
        // 不可达点的势能设为 0，避免后续计算 INF 参与运算
        for (int i = 1; i <= n; i++) {
            if (h[i] == INF) h[i] = 0;
        }
    }

public:
    MinCostMaxFlow(int n) : n(n), g(n + 1) {}

    // 加一条 u -> v 的边，容量为 cap，费用为 cost
    // 同时加反向边 v -> u，容量为 0，费用为 -cost，用于撤销流量
    void addEdge(int u, int v, ll cap, ll cost) {
        int id = edges.size();
        edges.push_back({u, v, cap, cost});
        edges.push_back({v, u, 0, -cost});
        g[u].push_back(id);
        g[v].push_back(id ^ 1);
    }

    // 返回 {最大流, 最小费用}
    // useSPFA 表示是否使用 SPFA 初始化势能
    pair<ll, ll> minCostMaxFlow(int s, int t, bool useSPFA = false) {
        ll maxFlow = 0;
        ll minCost = 0;

        vector<ll> h(n + 1, 0); // 势能，用来让 Dijkstra 可以处理残量网络中的负边
        vector<ll> dist(n + 1); // 最短路距离
        vector<int> pre(n + 1); // pre[v] 表示最短路上到达 v 的边编号

        if (useSPFA) {
            initPotentialBySPFA(s, h);
        }

        while (true) {
            fill(all(dist), INF);
            fill(all(pre), -1);
            priority_queue<pair<ll, int>, vector<pair<ll, int>>, greater<pair<ll, int>>> pq;
            dist[s] = 0;
            pq.push({0, s});
            // 在残量网络上跑 Dijkstra，找费用最短的增广路
            while (!pq.empty()) {
                auto [d, u] = pq.top();
                pq.pop();
                if (d != dist[u]) continue;
                for (int id : g[u]) {
                    Edge &e = edges[id];
                    if (e.cap <= 0) continue;
                    int v = e.to;
                    // 势能修正后的边权
                    // 如果势能正确，这个 newCost 应该是非负的
                    ll newCost = e.cost + h[u] - h[v];

                    if (dist[v] > dist[u] + newCost) {
                        dist[v] = dist[u] + newCost;
                        pre[v] = id;
                        pq.push({dist[v], v});
                    }
                }
            }
            // t 不可达，说明已经没有增广路
            if (pre[t] == -1) break;

            // 更新势能
            for (int i = 1; i <= n; i++) {
                if (dist[i] < INF) {
                    h[i] += dist[i];
                }
            }
            // 找这条增广路上的瓶颈容量
            ll addFlow = INF;
            for (int cur = t; cur != s; cur = edges[pre[cur]].from) {
                addFlow = min(addFlow, edges[pre[cur]].cap);
            }
            // 沿着最短费用路增广
            for (int cur = t; cur != s; cur = edges[pre[cur]].from) {
                int id = pre[cur];
                edges[id].cap -= addFlow;
                edges[id ^ 1].cap += addFlow;

                minCost += addFlow * edges[id].cost;
            }
            maxFlow += addFlow;
        }
        return {maxFlow, minCost};
    }
};
\end{minted}




\section*{字符串}
\addcontentsline{toc}{section}{字符串}


下面的字符串算法均默认采用 0-base实现 \par

\subsection*{KMP}
\addcontentsline{toc}{subsection}{KMP}


\begin{minted}{cpp}
 	  string s1;
    cin>>s1;
    const int n = s1.size();
		//这里采用0-base
    vector<int>next(n);
    for(int i = 1; i < n ;i ++){
      int len = next[i - 1];
      while(len && s1[i] != s1[len]){
        len = next[len - 1];
      }
      if(s1[i] == s1[len])
      next[i] = len + 1;
    }

    cout << n - next[n - 1] ; 
\end{minted}




\subsection*{manacher 求最长回文半径}
\addcontentsline{toc}{subsection}{manacher 求最长回文半径}


\begin{minted}{cpp}
  string s;
  cin >> s;
  const int n = s.size();
  vector<int> manacher(n);
  // 奇回文半径,manacher[cur] 表示以 cur 为中心的最大回文半径
  for (int cur = 0, L = 0, R = 0, k = 0; cur < n; cur++)
  {
    if (cur <= R)
      k = min(manacher[R + L - cur], R - cur + 1);
    else
      k = 1;
    while (cur + k < n && cur - k >= 0 && s[cur + k] == s[cur - k])
      k++;
    if (cur + k - 1 > R)//一定保证右端点最右
    	R = cur + k - 1, L = cur - k + 1;
    manacher[cur] = k;
  }
	fill(all(manacher),0);
  // 偶回文半径，manacher[cur] 表示以 {cur - 1, cur}为中心的最大回文半径
  for (int cur = 0, L = 0, R = 0, k = 0; cur < n; cur++)
  {
    if (cur <= R)
    k = min(manacher[R + L - cur + 1], R - cur + 1);
    else
    k = 0; // 注意这里要写成0，因为第一个位置初始不合法
    while (cur + k < n && cur - 1 - k >= 0 && s[cur + k] == s[cur - 1 - k])
    k++;
    
    if (cur + k - 1 > R)//一定保证右端点最右
    	R = cur + k - 1, L = cur - 1 - k + 1;
    manacher[cur] = k;
  }
\end{minted}


\subsection*{Z函数求所有后缀与原字符串的LCP}
\addcontentsline{toc}{subsection}{Z函数求所有后缀与原字符串的LCP}


\begin{minted}{cpp}
 // Z函数:求字符串每一个后缀和原字符串LCP
  auto Z = [&](string s) -> vector<int>
  {
    int n = s.size();
    vector<int> Z(n);

    // 默认第一个位置是全匹配
    Z[0] = n;
    for (int cur = 1, l = 0, r = 0, k = 0; cur < n; cur++)
    {
      if (cur <= r)
        k = min(Z[cur - l], r - cur + 1); // 复用维护的匹配区间信息
      else
        k = 0;
      // 暴力扩展
      while (cur + k < n && s[cur + k] == s[k])
        k++;

      // 当且仅当新匹配区间右端点更远的时候更新匹配区间
      if (cur + k - 1 > r)
        l = cur, r = cur + k - 1;
      // 记录Z函数的值
      Z[cur] = k;
    }
    return Z;
  };
\end{minted}




\section*{动态规划}
\addcontentsline{toc}{section}{动态规划}


\subsection*{SOSDP}
\addcontentsline{toc}{subsection}{SOSDP}


子集和DP ：给定 $$ val[mask]$$，求 $$ dp[mask] = \sum{  val[sub] [sub ⊆ mask ]}$$  。 \par

\begin{minted}{cpp}
// 子集和 SOS DP
// dp[mask] = sum(val[sub])，其中 sub 是 mask 的子集
// 复杂度 O(n * 2^n),适用于n <= 20甚至25
  for (int bit = 0; bit < n; bit++)
  {
    for (int mask = 0; mask < (1 << n); mask++)
    {
      // 如果 mask 含有 bit，那么它的子集可以分成两类：
      // 1. 含 bit 的子集，已经在 dp[mask] 里
      // 2. 不含 bit 的子集，来自 dp[mask ^ (1 << bit)]
      if ((mask >> bit) & 1)//把这里的值等于1改成0即为超集和DP
      {
        dp[mask] += dp[mask ^ (1 << bit)];
      }
    }
  }
\end{minted}


\subsection*{树上背包}
\addcontentsline{toc}{subsection}{树上背包}


\begin{minted}{cpp}
  vector<vector<ll>> dp1(n + 1,vector<ll>(n + 1)), dp2(n + 1,vector<ll>(n + 1));
  vector<ll> sz(n + 1);
  ll ans = 0 ;
  [&](this auto && dfs,int u, int p) -> void
  {
    sz[u] = 1;
    dp1[u][1] = 1;
    for (int v : adj[u])
    {
      if (v == p)
        continue;
        dfs(v,u);
      vector<ll>ndp1 =dp1[u];
      vector<ll>ndp2 =dp2[u];
      //注意正是这里按照sz合并的优化，让每一个点对只会在自己的LCA处被枚举到一次，故复杂度是O(n^2)的
      for(int i = 1;i <= sz[v] ;i ++){
        ndp1[i + 1] += dp1[v][i];
        ndp2[i + 1] += dp2[v][i];
      }
      //假如不开滚动数组，需要和01背包一样倒序枚举内层
      for (int i = 1; i <= sz[v]; i++)
      {
        for (int j = 1; j <= sz[u]; j++)
        {
          ndp2[i + j] += dp1[u][j] * dp1[v][i];
          if(i + j == k){
            ans += dp1[u][j] * dp2[v][i] + dp2[u][j] * dp1[v][i];
          }
        }
      }
      sz[u] += sz[v];
      swap(dp1[u],ndp1);
      swap(dp2[u],ndp2);
    }
  }(1,0);
\end{minted}


\subsection*{多重背包（二进制拆分）}
\addcontentsline{toc}{subsection}{多重背包（二进制拆分）}


\begin{minted}{cpp}
 		vector<int> t(n+1), c(n+1), p(n+1);
    for (int i = 1; i <= n; ++i) cin >> t[i] >> c[i] >> p[i];
    // dp[cap] = 最大美学值
    vector<int> dp(T+1, 0);

    for (int i = 1; i <= n; ++i) {
        // 跳过必定用不上的物品（可选，主要是小优化）
        if (t[i] > T) continue;

        if (p[i] == 0) {
            // 完全背包：正序
            for (int cap = t[i]; cap <= T; ++cap)
                dp[cap] = max(dp[cap], dp[cap - t[i]] + c[i]);
        } else {
            // 多重背包：二进制拆分 => 0/1 物品，倒序
            int x = p[i];
            int k = 1;
            while (x > 0) {
                int num = min(k, x);
                int cost = t[i] * num, val = c[i] * num;
                if (cost <= T) {
                    for (int cap = T; cap >= cost; --cap)
                        dp[cap] = max(dp[cap], dp[cap - cost] + val);
                }
                x -= num;
                k <<= 1;
            }
        }
    }
\end{minted}




\section*{几何}
\addcontentsline{toc}{section}{几何}


\subsection*{经典结论与技巧}
\addcontentsline{toc}{subsection}{经典结论与技巧}


1. 对于平面上的点 \texttt{(X, Y)}，如果要\textbf{最大化}： \par

\begin{minted}{cpp}
   |Y1 - Y2| / |X1 - X2|
\end{minted}


只需要 **按照 \texttt{X} 排序**，然后检查相邻点。 \par

原因是：如果最优点对在 \texttt{X} 排序中间隔了别的点，那么中间那个点一定能和两端之一组成一个不劣的点对。 \par

\section*{其它技巧}
\addcontentsline{toc}{section}{其它技巧}


\subsection*{普通莫队}
\addcontentsline{toc}{subsection}{普通莫队}


\begin{minted}{cpp}
    int n,m,k;
    cin>>n>>m>>k;
    vector<int>a(n + 1);
    for(int i =1; i<=n ;i ++)cin>>a[i];
    vector<tuple<int,int,int>>todo;
    todo.reserve(m);
    vector<int>ans(m + 1);
    for(int i= 1; i <= m;i ++)
    {
      int l,r;cin>>l>>r;
      todo.emplace_back(l,r,i);
    }
    const int M = sqrt(n);
    sort(all(todo),[&](tuple<int,int,int>a,tuple<int,int,int>b){
      auto[l_a,r_a,_] = a;
      auto[l_b,r_b,_] = b;
      int id_a = (l_a - 1)/M + 1;
      int id_b = (l_b - 1)/M + 1;
      if(id_a != id_b){
        return id_a < id_b;
      }
      if(id_a & 1){
        return r_a < r_b;
      }else{
        return r_a > r_b;
      }
    });
    ll sum = 0;
    int wl = 1,wr = 0;

    vector<ll>cnt(k + 1,0);
    auto del = [&](int p){
      sum -= (2*cnt[a[p]] - 1);
      cnt[a[p]]--;
    };
    auto add = [&](int p){
      sum += (2*cnt[a[p]] + 1);
      cnt[a[p]]++;
    };
    for(auto&[l,r,id] : todo)
    {
      while(wl > l){
        add(--wl);
      }
      while(wl < l){
        del(wl++);
      }
      while(wr > r){
        del(wr--);
      }
      while(wr < r){
        add(++wr);
      }
      
      ans[id] = sum;
    }
    for(int i = 1;i <= m;i ++)
    cout<<ans[i]<<"\n";
\end{minted}


\subsection*{带修莫队}
\addcontentsline{toc}{subsection}{带修莫队}


\begin{minted}{cpp}
struct Event { char typ; int x, y; };            // Q: x=L,y=R;  R: x=P,y=C(original)
struct Query  { int l, r, t, id; };              // t = updates done before this query
struct Update { int p, from, to; };              // color change at position p

int main() {
  cin.tie(0)->sync_with_stdio(0);

  int n, m; 
  cin >> n >> m;
  vector<int> A(n + 1);
  for (int i = 1; i <= n; i++) cin >> A[i];

  vector<Event> event; event.reserve(m);
  vector<int> color; color.reserve(n + m);
  for (int i = 1; i <= n; i++) color.push_back(A[i]);

  for (int i = 0; i < m; i++) {
    char op; int x, y;
    cin >> op >> x >> y;            
    event.push_back({op, x, y});
    if (op == 'R' || op == 'C') color.push_back(y); 
  }
  sort(all(color));
  color.erase(unique(all(color)), color.end());
  auto get_id = [&](int v)->int {
    return int(lower_bound(all(color), v) - color.begin()) + 1; 
  };
  for (int i = 1; i <= n; i++) A[i] = get_id(A[i]);

  vector<Query> query; query.reserve(m);
  vector<Update> update; update.reserve(m);
  vector<int> a_cur = A;   
  int time = 0, idx = 0;
  for (auto [tp, x, y] : event) {
    if (tp == 'Q') {
      query.push_back({x, y, time, ++idx});
    } else { 
      int p = x, to = get_id(y);
      update.push_back({p, a_cur[p], to});
      a_cur[p] = to;
      ++time;
    }
  }
  const int B = max(1, (int)pow(n, 2.0/3.0));
  auto id = [&](int x){ return x / B; };
  sort(all(query), [&](const Query& A, const Query& Bq){
    int ab = id(A.l), bb = id(Bq.l);
    if (ab != bb) return ab < bb;
    int ar = id(A.r), br = id(Bq.r);
    if (ar != br) return (ab & 1) ? (ar > br) : (ar < br);
    return ( (ar & 1) ? (A.t > Bq.t) : (A.t < Bq.t) );
  });
  int K = (int)color.size() + 5;
  vector<int> cnt(K, 0);
  vector<int> ans(idx + 1, 0);
  vector<int> a = A;    
  int l = 1, r = 0, curT = 0, kind = 0;

  auto add = [&](int pos){
    pos = a[pos];
    if (++cnt[pos] == 1) ++kind;
  };
  auto del = [&](int pos){
    pos = a[pos];
    if (--cnt[pos] == 0) --kind;
  };
  auto modify = [&](){      
    auto& [p,from,to] = update[curT];
    if (l <= p && p <= r) {
      if (--cnt[from] == 0) --kind;
      if (++cnt[to]   == 1) ++kind;
    }
    a[p] = to;
    ++curT;
  };
  auto recover = [&](){            
    auto&[p,from,to] = update[curT - 1];
    if (l <= p && p <= r) {
      if (--cnt[to]   == 0) --kind;
      if (++cnt[from] == 1) ++kind;
    }
    a[p] = from;
    --curT;
  };

  for (auto&[L,R,T,id] : query) {
    while (curT < T) modify();
    while (curT > T) recover();
    while (r < R) add(++r);
    while (l > L) add(--l);
    while (r > R) del(r--);
    while (l < L) del(l++);
    ans[id] = kind;
  }

  for (int i = 1; i <= idx; i++) cout << ans[i] << "\n";
  return 0;
}
\end{minted}




\end{multicols*}
\end{document}
